\documentclass[letterpaper]{article}
% DO NOT CHANGE THIS
\usepackage{aaai2026} % DO NOT CHANGE THIS
\usepackage{times} % DO NOT CHANGE THIS
\usepackage{helvet} % DO NOT CHANGE THIS
\usepackage{courier} % DO NOT CHANGE THIS
\usepackage[hyphens]{url} % DO NOT CHANGE THIS
\usepackage{graphicx} % DO NOT CHANGE THIS
\urlstyle{rm} % DO NOT CHANGE THIS
\def\UrlFont{\rm} % DO NOT CHANGE THIS
\usepackage{graphicx} % DO NOT CHANGE THIS
\usepackage{natbib} % DO NOT CHANGE THIS
\usepackage{caption} % DO NOT CHANGE THIS
\frenchspacing % DO NOT CHANGE THIS
\setlength{\pdfpagewidth}{8.5in} % DO NOT CHANGE THIS
\setlength{\pdfpageheight}{11in} % DO NOT CHANGE THIS
%
% Keep the \pdfinfo as shown here. There’s no need
% for you to add the /Title and /Author tags.
\pdfinfo{
/TemplateVersion (2026.1)
}

%File: anonymous-submission-latex-2026.tex


% \columnsep -- This command may not be used
% \newpage -- No page breaks of any kind may be used for the final version of your paper
% \pagebreak -- No page breaks of any kind may be used for the final version of your paperr
% \pagestyle -- This command may not be used
% \tiny -- This is not an acceptable font size.
% \vspace{- -- No negative value may be used in proximity of a caption, figure, table, section, subsection, subsubsection, or reference
% \vskip{- -- No negative value may be used to alter spacing above or below a caption, figure, table, section, subsection, subsubsection, or reference

% ========================

\usepackage{multirow}%
\usepackage{makecell}%
\usepackage{amsmath,amssymb,amsfonts}%
\usepackage{mathrsfs}%
%\usepackage[title]{appendix}%
\usepackage{appendix}%
\usepackage{xcolor}%
\usepackage{textcomp}%
\usepackage{manyfoot}%
\usepackage{booktabs}%
%\usepackage{algorithm}%
%\usepackage{algorithmicx}%
\usepackage{algpseudocode}%
\usepackage[ruled,linesnumbered]{algorithm2e}%
\usepackage{listings}%
\usepackage{enumitem}

% \usepackage{hyperref}  % not allowed
\usepackage{cleveref}
\usepackage{subcaption}
\usepackage{booktabs}
\usepackage{placeins}
\usepackage{calc}  % For inkscape's \svgscale

\usepackage{mdframed}
\usepackage{amsthm}
\usepackage{thmtools, thm-restate}
\usepackage{tikz-cd}

\graphicspath{{figures},{figures/illustrations}}

\usepackage{longtable}

%\DeclareEmphSequence{\bfseries}  % Set \emph to bold
\DeclareEmphSequence{}  % disable \emph

\newcommand{\T}{^T}
\newcommand{\x}{\mathbf{x}}
\newcommand{\y}{\mathbf{y}}
\newcommand{\X}{\mathcal{X}}
\newcommand{\Y}{\mathcal{Y}}
\newcommand{\M}{\mathcal{M}}
\newcommand{\PP}{\mathcal{P}}
\newcommand{\Pc}{\mathcal{P}_c}
% \newcommand{\Pr}{\mathcal{P}_r}
\newcommand{\A}{\mathcal{A}}
\newcommand{\B}{\mathcal{B}}
\newcommand{\C}{\mathcal{C}}
\newcommand{\I}{\mathcal{I}}
\newcommand{\s}{\mathbf{s}}
\newcommand{\R}{\mathbb{R}}
\newcommand{\N}{\mathbb{N}}

%% as per the requirement new theorem styles can be included as shown below
%\theoremstyle{thmstyleone}%
\newtheorem{theorem}{Theorem}[section]
\newtheorem{proposition}[theorem]{Proposition}% 
\newtheorem{corollary}{Corollary}[theorem]
\newtheorem{lemma}{Lemma}

%\theoremstyle{thmstyletwo}%
\newtheorem{example}{Example}%
\newtheorem{remark}{Remark}%

%\theoremstyle{thmstylethree}%
\newtheorem{definition}{Definition}%
\newtheorem{assumption}{Assumption}%

\declaretheoremstyle[
    headfont=\bfseries, 
    bodyfont=\normalfont\itshape,
    headpunct={},
    spacebelow=\parsep,
    spaceabove=\parsep,
    mdframed={
        innertopmargin=6pt,
        innerbottommargin=6pt, 
        skipabove=\parsep, 
        skipbelow=\parsep} 
]{framedstyle}

\declaretheorem[
    style=framedstyle,
    name=Theorem
]{thm}

%\raggedbottom

% Add appendix to cleverref
% \AddToHook{cmd/appendix/before}{%
%     \crefalias{section}{appendix}%
%     \crefalias{subsection}{appendix}
% }
% ========================
%
% These are recommended to typeset algorithms but not required. See the subsubsection on algorithms. Remove them if you don't have algorithms in your paper.
% \usepackage{algorithm}
% \usepackage{algorithmic}

%
% These are are recommended to typeset listings but not required. See the subsubsection on listing. Remove this block if you don't have listings in your paper.
\usepackage{newfloat}
% \usepackage{listings}
\DeclareCaptionStyle{ruled}{labelfont=normalfont,labelsep=colon,strut=off} % DO NOT CHANGE THIS

% Title

% Your title must be in mixed case, not sentence case.
% That means all verbs (including short verbs like be, is, using,and go),
% nouns, adverbs, adjectives should be capitalized, including both words in hyphenated terms, while
% articles, conjunctions, and prepositions are lower case unless they
% directly follow a colon or long dash
\title{Oxytrees: Model Trees for Bipartite Learning}
\author{
    Pedro Ilídio\textsuperscript{\rm 1,\rm 2},
    Felipe Kenji Nakano\textsuperscript{\rm 1,\rm 2},
    Alireza Gharahighehi\textsuperscript{\rm 1,\rm 2},\\
    Robbe D'Hondt\textsuperscript{\rm 1,\rm 2},
    Ricardo Cerri\textsuperscript{\rm 3},
    Celine Vens\textsuperscript{\rm 1,\rm 2}
}

\affiliations{
    \textsuperscript{\rm 1}
        KU Leuven, Campus KULAK,
        Dept. of Public Health and Primary Care
        Etienne Sabbelaan 53, Kortrijk, 8500, Belgium\\
        
    \textsuperscript{\rm 2}
        Itec, imec research group at KU Leuven,
        Etienne Sabbelaan 51, Kortrijk, 8500, Belgium\\
        
    \textsuperscript{\rm 3}
        Universidade de São Paulo,
        Instituto de Ciências Matemáticas e de Computação
        Av. Trab. São Carlense, 13566-590, São Carlos, São Paulo, Brazil\\
        
    \{%
    pedro.ilidio,
    felipekenji.nakano,
    alireza.gharahighehi,
    robbe.dhondt,
    celine.vens%
    \}@kuleuven.be, 
    cerri@icmc.usp.br
    
    % pedro.ilidio@kuleuven.be
    % felipekenji.nakano@kuleuven.be
    % alireza.gharahighehi@kuleuven.be
    % robbe.dhondt@kuleuven.be
    % cerri@icmc.usp.br
    % celine.vens@kuleuven.be
}


%\setcounter{secnumdepth}{0} %May be changed to 1 or 2 if section numbers are desired.
\setcounter{secnumdepth}{2} %May be changed to 1 or 2 if section numbers are desired.


% Show lengths useful for plotting configuration 
% \usepackage{layouts}
% \begin{document}
% {
%     \color{red}
%     \par\noindent TEXT\_WIDTH = \printinunitsof{in}\prntlen{\textwidth}
%     \par\noindent COLUMN\_WIDTH = \printinunitsof{in}\prntlen{\columnwidth}
% }
% \end{document}

% TEXT WIDTH = 7.00137in
% COLUMN WIDTH = 3.31314in

\begin{document}

\maketitle


\begin{abstract}
Bipartite learning is a machine learning task that aims to predict interactions between pairs of instances. It has been applied to various domains, including drug-target interactions, RNA-disease associations, and regulatory network inference. Despite being widely investigated, current methods still present drawbacks, as they are often designed for a specific application and thus do not generalize to other problems or present scalability issues.
To address these challenges, we propose Oxytrees: proxy-based biclustering model trees. Oxytrees compress the interaction matrix into row- and column-wise proxy matrices, significantly reducing training time without compromising predictive performance.
We also propose a new leaf-assignment algorithm that significantly reduces the time taken for prediction.
Finally, Oxytrees employ linear models using the Kronecker product kernel in their leaves, resulting in shallower trees and thus even faster training. Using 15 datasets, we compared the predictive performance of ensembles of Oxytrees with that of the current state-of-the-art.
We achieved up to 30-fold improvement in training times compared to state-of-the-art biclustering forests, while demonstrating competitive or superior performance in most evaluation settings, particularly in the inductive setting.
Finally, we provide an intuitive Python API to access all datasets, methods and evaluation measures used in this work, thus enabling reproducible research in this field.
\end{abstract}

\begin{links}
\link{Code}{https://github.com/pedroilidio/oxytrees2025}
%\link{Python package}{https://github.com/pedroilidio/bipartite_learn}
\link{Python package}{https://bipartite-learn.readthedocs.io}
%\link{Datasets}{https://aaai.org/example/datasets}
\link{Extended version}{https://arxiv.org/abs/2511.12713}
\end{links}


\section{Introduction}
\label{sec:introduction}

Over the last two decades, technological advancements have led to the daily generation of enormous volumes of data, encompassing social media interactions, e-commerce transactions, IoT sensor outputs, and genomic sequences.
Such data are frequently presented in complex structures, which challenge the application of traditional machine learning algorithms.
In this work, we focus on the problem of bipartite learning, a machine learning task that involves predicting interactions within a network where two distinct types of instances are interconnected.
%
More specifically, bipartite learning is the task of modeling a function $(x_1,\; x_2) \mapsto y$ mapping a pair of objects $x_1$ and $x_2$ of different types to an output $y$ characterizing their interaction. Often, these datasets can be represented by two feature matrices, $X_1$ and $X_2$, which describe instances in each dimension, and an interaction matrix, $Y$, that holds their corresponding output~values. % (\cref{fig:bipartite dataset}).


Bipartite learning has been applied to \emph{several contexts}.
For instance, in bioinformatics, deep learning methods are employed for drug-target interaction \cite{huang_moltrans_2020,lin2023comprehensive, bagherian2021machine,liu2024ssldti}, miRNA-disease association \cite{tian2024mgcnss} and compound-protein interaction prediction \cite{tsubaki2019compound}.
Similarly, recommender systems can be seen as a user-item interaction problem \cite{aggarwal2016recommender}.

Regardless of the application, validating all possible interactions in $Y$ is unfeasible, e.g., it is impractical to verify all possible user-item or drug-target interactions. Consequently, many interactions are not directly observed but instead predicted using computational models trained on available data. It is often assumed that some of the training interactions labeled as negative are, in fact, unconfirmed positives, a setting commonly referred to as \emph{positive-unlabeled} (PU) learning~\cite{bekker2020learning}.
Many proposed methods can only predict interactions among instances present in the training set, i.e., they are not applicable to new instances~\cite{park_flaws_2012,schrynemackers_protocols_2013,pahikkala_toward_2015}, and hence, are not inductive~\cite{chapelle_semi-supervised_2006}. This limitation presents a significant challenge in the context of recommender systems, where it is referred to as the cold-start problem~\cite{gharahighehi_addressing_2022}. In the context of drug-protein interaction prediction, it is crucial to generalize to new disease proteins.
For that purpose, state-of-the-art techniques rely on deep learning to model the complex structure of the inputs (e.g., proteins as sequences, molecules as graphs)~\cite{huang_moltrans_2020,dehghan_tripletmultidti_2023}. As a result, they can only solve a specific type of interaction, severely compromising interpretability, and are computationally expensive, requiring large amounts of training data.

In this work, we investigate decision tree-based methods, due to their applicability to any type of interaction problem,
interpretability, capacity to handle problems with few training data, and state-of-the-art performance on tabular input~\cite{grinsztajn_why_2022,costa2023recent,blockeel2023decision}. Although promising, the current tree-based approaches to bipartite learning also pose significant scalability limitations. Furthermore, these decision trees fail to accurately represent sparse regions of the training space, since each leaf is a very simple model (the average label) of a small localized region.
% 
In this paper, we propose a novel method, named \emph{Oxytrees}, as a solution to these limitations.


The core idea of Oxytrees is to aggregate the impurity of one of the dimensions of $Y$ into small matrices we call \textit{proxies}. The proxies are reused when evaluating multiple split candidates, aggregating only the remaining dimension when calculating the impurity of each candidate.
Additionally, Oxytrees introduce 
a \emph{faster leaf-assignment procedure}, for quicker inference of all interactions between large batches of instances; as well as
\emph{model trees for bipartite learning}.
\emph{Model trees}~\cite{quinlan1992learning,landwehr_logistic_2005,costa2023recent} return a (simple) function in the leaf nodes, rather than a constant.
This way, instance pairs arriving in the same leaf can still get very different predictions, allowing to more easily model complex relationships, and at the same time potentially allowing smaller tree sizes (with larger leaves\footnote{\textit{Oxytree} is the popular name of a genus of trees (\textit{Paulownia sp.}) that are not too tall, grow very fast, and have very large leaves.}), hence reducing the tree construction time even further.
Leaf models also enable trees to model functions that extend beyond the dense regions of the training set, where traditional decision trees struggle.

Two \emph{preliminary investigations} of our research groups relate to this study. \citet{alves2022twostep} use a single tree with gradient boosting models in the leaves with large label imbalance. The complex leaf models amplified the scalability issues.
We later~\cite{ilidio_fast_2024} introduced a preliminary version of the node-splitting procedure used by Oxytrees. In both studies, only a limited set of experiments was evaluated, and the performance improvements remained unclear.

For the present study, we built ensembles of Oxytrees and compared them to state-of-the-art techniques in bipartite learning. We selected methods applicable to a wide range of inductive interaction prediction tasks and compared them across 15 datasets.
Furthermore, we investigate the PU setting in detail by randomly masking 25\%, 50\%, and 75\% of the positive interactions.
%
To the best of our knowledge, this represents the most extensive comparative study to date in the context of inductive bipartite learning, which helps delineate the state-of-the-art of this emerging paradigm.
%
To assist future studies and allow reproducible research, we make all our algorithm implementations and preprocessed datasets \emph{publicly available},\footnote{\url{https://anonymous.4open.science/r/oxytrees2025-004C}}
using an accessible and intuitive Python API based on Scikit-Learn~\cite{pedregosa_scikit-learn_2011}.


\section{Formal Problem Definition}
\label{sec:background}


Consider two independent feature spaces $\X_1$ and $\X_2$, and an output space $\Y$.
We address the problem of modeling a function $f \colon \X_1 \times \X_2 \to \Y$ given only finite samples $X_1 \in \X_1^{n_1}$, $X_2 \in \X_2^{n_2}$, and $Y\in \Y^n$, with $Y$ representing known outputs of $f$ corresponding to dyads in $X_1 \times X_2$.
We refer to this paradigm as \emph{bipartite learning}.
We call $\X_1$ and $\X_2$ the two \emph{domains} or \emph{dimensions} of the problem.

In our specific learning setting, we focus on tasks with structured input with $\X_1 \subseteq \R^{m_1}$ and $\X_2 \subseteq \R^{m_2}$. Therefore, $X_1$ and $X_2$ can be represented as matrices $X_1=(X_1^{ij})$ and $X_2=(X_2^{ij})$. We use indices as superscripts and reserve the subscripts for descriptors of the matrix or vector as a whole.
%
We consider a single output, so $\Y \subseteq \R$ and an interaction matrix $Y$ can be built to hold the annotated interactions as $Y^{ij} = f(X_1^i,\; X_2^j)$.
%
We specifically focus on binary classification and therefore $\Y=\{0,\;1\}$. Furthermore, we assume the PU setting~\cite{bekker2020learning}, meaning that positive annotations are deemed confirmed, but negative annotations can be either ground-truth negatives or unannotated positives.


\section{Related Work}
\label{sec:related work}

A handful of terms have been used in the literature to describe similar learning problems,
such as interaction prediction~\cite{schrynemackers_classifying_2015,pliakos_network_2019,chen_machine_2018,bagherian_machine_2020}, link prediction~\cite{lu_link_2011,zhou_progresses_2021}, dyadic prediction~\cite{menon_log-linear_2010,pahikkala_two-step_2014,jin_multitask_2017}, and network inference~\cite{park_flaws_2012,pahikkala_toward_2015,schrynemackers_protocols_2013,pliakos_network_2019}.
Link prediction, dyadic prediction, and network inference do not specify the existence of feature spaces ($X_1$ and $X_2$ can be derived from $Y$). Furthermore, they allow $X_1 = X_2$, not resulting in a bipartite network, as is also the case for interaction prediction.
% they are inductive, as opposed to transductive

To build inductive models for bipartite learning,
two common approaches reorganize the data to enable traditional algorithms. We call them \emph{data-based} adaptations, and they are classified as global and local~\cite{schrynemackers_classifying_2015}.
Global methods transform $X_1$ and $X_2$ into a single input space, whereas local methods build separate models for each domain~\cite{schrynemackers_classifying_2015}.

Unlike data-based approaches, some methods propose or modify machine learning estimators on a fundamental level to take maximum advantage of the bipartite learning setting.
We call them estimator-based adaptations.
\emph{Estimator-based} methods from the recent literature are mainly based on deep learning, matrix factorization or decision trees.
\emph{Deep learning} methods provide end-to-end solutions in which data are processed in their raw format. For example, graphs (compounds) and sequences (proteins), together with their known interactions, are directly used as input to the neural networks in drug-target bipartite learning.
\emph{Matrix factorization} methods focus on decomposing the interaction matrix in a lower-dimensional latent representation, which is then used to infer interactions~\cite{guo2024rise,chen2022review}. Neighborhood-Regularized Logistic Matrix Factorization (NRLMF) carries neighborhoods to the latent space to enable inductiveness~\cite{liu_neighborhood_2016,liu_lpi-nrlmf_2017,liu_predicting_2020}.
%
\emph{Decision tree-based methods}~\cite{costa2023recent,blockeel2023decision} are relatively new in bipartite learning. Predictive Bi-Clustering Trees (PBCT)~\cite{pliakos_global_2018} are an extension of Predictive Clustering Trees \cite{blockeel1998top} that allows splits using features from both $X_1$ and $X_2$.
In this way, a biclustering (i.e. two-dimensional partitioning) of the interaction matrix is obtained.
PBCT was then extended to ensembles~\cite{pliakos_network_2019}, and further combined with a pre-training step that imputes the unreliable negative annotations using NRLMF.
The resulting method was named BICTR, and showed state-of-the-art performance among the methods that are applicable to all bipartite learning problems~\cite{pliakos_drug-target_2020}.
Nonetheless, the method was only applied to the drug-target setting, and has scalability limitations both in induction and inference time. 


\section{The Proposed Method: Oxytrees}
\label{sec:oxytrees}

This study introduces Oxytrees, proxy-based model trees for bipartite learning. It combines a refined training algorithm with a new batch inference procedure and leaf-specific models to generate outputs.


\subsection{Requirements for the Impurity Function}
\label{sec:oxytrees impurity}

Let $Y_\text{node}$ be a subset of
%$n_\text{node}$ dyads of
$Y$ reaching a given tree node.
Decision trees use an impurity function $I(Y_\text{node})$ to select a split in each tree node~\cite{breiman_classification_1984,blockeel1998top}.
%
Oxytrees require (see theorem C2) that $I(Y_\text{node})$ can be expressed in terms of two arbitrary functions $\mu$ and $\rho$:
%Oxytrees only require $I(Y)$ to be invariant under permutations of elements in $Y_\text{node}$.
%Oxytrees only require $I(Y)$ to be invariant to the order in which dyads are presented during training (i.e. dyad permutation invariance).
%This requirement is implicitly assumed for any method that considers each dyad as an independent instance, which corresponds to the vast majority of current methods. Dyad permutation invariance holds if and only if $I(Y_\text{node})$ can be expressed in terms of two arbitrary functions $\mu$ and $\rho$ (see Theorems 1 and 2 in \cref{sec:proof}):
\begin{equation}
    \label{eq:i constraint}
    \textstyle
    I_\text{Oxytrees} (Y_\text{node}) = \rho(\sum_{ij}\mu(Y_\text{node}^{ij}))
    \text{.}
\end{equation}
%
In this study, we specifically used the variance of the vectorized interaction matrices: $I_\text{Oxytrees} (Y_\text{node}) = \sigma^2(\text{vec}(Y_\text{node}))$. More explicitly:
%
\begin{equation}
    \label{eq:i bgso}
    I_\text{Oxytrees} (Y_\text{node})
    %= \sigma^2(\text{vec}(Y_\text{node}))
        %= \frac{\sum_{ij}(Y^{ij})^2}{|Y||Y^T|} - \left(\frac{\sum_{ij}Y^{ij}}{|Y||Y^T|}\right)^2
        %= \frac{\sum_{ij}(Y^{ij}_\text{node})^2}{n_\text{node}} - \left(\frac{\sum_{ij}Y^{ij}_\text{node}}{n_\text{node}}\right)^2
    = \frac{\sum_{ij}(Y^{ij}_\text{node})^2}{n_\text{node}} - \left(\frac{\sum_{ij}Y^{ij}_\text{node}}{n_\text{node}}\right)^2
        %= \frac{\sum_{ij}(Y^{ij}_\text{node})^2}{n_\text{node}} - \frac{(\sum_{ij}Y^{ij}_\text{node})^2}{n_\text{node}^2}
    \text{,}
\end{equation}
which corresponds to
\begin{align}
    \mu(z) = \begin{bmatrix} 1, & z, & z^2 \end{bmatrix}
    &&
    \rho(z_1, z_2, z_3) = \frac{z_3}{z_1} - \frac{z_2^2}{z_1^2}
    \text{.}
\end{align}

The current state-of-the-art BICTR~\cite{pliakos_drug-target_2020} considers columns or rows of $Y_\text{node}$ as multiple outputs, using $I(Y_\text{node})$ as the average of the variances of each column or row. This is incompatible with \cref{eq:i constraint} (see corollary C2.1) and therefore prohibits the proposed training optimization.

More intuitively, \cref{eq:i constraint} is equivalent to requiring $I(Y)$ to be invariant to the order in which dyads are presented during training (see theorem C1).
This is a standard assumption, the only exceptions being BICTR-related methods and local methods~\cite{schrynemackers_protocols_2013}.


\subsection{Using Proxies to Speed Up Training}
\label{sec:oxytrees training}

The impurity format in \cref{eq:i constraint} enables a significant improvement in computational performance.
%
Oxytrees start the split search in a node by building proxy matrices $\tilde Y_1$ and $\tilde Y_2$ that aggregate row- and column-wise impurities, respectively:
%
\begin{align}
    \label{eq:proxies general}
    \textstyle
    \tilde Y_1^i = \sum_{j} \mu(Y_\text{node}^{ij})
    &&
    \textstyle
    \tilde Y_2^j = \sum_{i} \mu(Y_\text{node}^{ij})
    \text{.}
\end{align}
%$\tilde Y_1$ and $\tilde Y_2$ defined by
%Notice from \cref{eq:i bgso} that $I_\text{BGSO}$ 
Oxytrees then use $\tilde Y_1$ to evaluate all the horizontal split candidates in the node (that partition the rows of $Y_\text{node}$), and $\tilde Y_2$ to evaluate all the vertical split candidates in the node (that partition the columns).
Finally, the split $s(Y_\text{node}) = \{Y_\text{A},\; Y_\text{B}\}$ with the largest decrease in impurity $\Delta I$ is chosen:
\begin{equation}
    \label{eq:delta i}
    \textstyle
    \Delta I =
        \frac{1}{n}(
            n_\text{node}\;I(Y_\text{node})
            - n_\text{A}\;I(Y_\text{A})
            - n_\text{B}\;I(Y_\text{B})
        )
        \text{.}
\end{equation}
%\footnote{It is usual to also include a $\frac{1}{n_\text{root}}$ factor in the $\Delta I$ definition~\cite{breiman_classification_1984}. We choose to omit it, since it does not affect the split search at each node. It is used mostly for impurity-based stop criteria, feature importances and pruning, which are out of the scope of this study.}

The reasoning behind the computational improvement is as follows. A large number of different candidate splits is usually evaluated in each node. For a given horizontal (resp. vertical) split, it is much more efficient to calculate $\Delta I$ from the precomputed $\tilde Y_1$ (resp. $\tilde Y_2$) than from the original $Y$, as this bypasses aggregation across columns (resp. rows).
%due to the much smaller size of $\tilde Y_1$ and $\tilde Y_2$.
As a consequence, the overall split search is faster using $\tilde Y_1$ and $\tilde Y_2$ for split evaluation, even if considering the extra time needed to build the proxy matrices.
%
Precisely, building an Oxytree has
the following computational complexity
\begin{equation}
    \label{eq:complexity oxytree}
    %\text{\ref{alg:dt build}}_\text{Oxytree}(n, \;  m) \in \Theta(n^2 (\log \; n + m))
    \text{BuildOxytree} \in \Theta(n_1^2 (\log (n_1) + m))
    \text{,}
\end{equation}
%
in which $m$ represents the number of features chosen to be evaluated at each node and $n_1 \propto n_2$ is the number of instances in each dimension. For PBCT we have
\begin{equation}
    \label{eq:complexity pbct}
    %\text{\ref{alg:dt build}}_\text{Oxytree}(n, \;  m) \in \Theta(n^2 (\log \; n + m))
    \text{BuildPBCT} \in \Theta(m n_1^2 \log (n_1))
    \text{.}
\end{equation}
As a result, Oxytrees reduce the training time by a factor of $\log (n_1)$  if $m \in \Omega(\log (n_1))$ and by a factor of $m$ otherwise, in comparison to PBCT-based methods (including the state-of-the-art, BICTR). A detailed derivation of \cref{eq:complexity oxytree} and \cref{eq:complexity pbct} is presented in
% \cref{sec:complexity analysis},
appendix D1,
while an empirical confirmation experiment is explored in \cref{sec:empirical complexity}.
Please refer to algorithms
%\ref{alg:find best split bgso} and \ref{alg:find best split gmo}
FindBestSplitOxytree and FindBestSplitBICTR
in \cref{sec:additional algorithms}
for pseudo-codes on the split search procedures.


\subsection{Efficient Batch Inference Procedure}
\label{sec:batch inference}

%We introduce a scalable procedure to predict interactions between large batches of instances, suited for biclustering decision trees in general. We now describe this procedure and briefly discuss the resulting improvement in computational complexity.

Consider a set of test row instances $X_\text{1, test}$ and a set of test column instances $X_\text{2, test}$.
The current state-of-the-art algorithm~\cite{pliakos_global_2018} traverses the tree once for each pair in $X_\text{1, test} \times X_\text{2, test}$ to make predictions (see \cref{fig:train test split}). Instead, our inference procedure receives the two entire feature matrices and, at each tree node, uses the split rule to partition either $X_\text{1, test}$ or $X_\text{2, test}$, depending on whether the rule represents a horizontal or vertical split. Each partition is passed to the corresponding child node, while the other, nonpartitioned, $X$ matrix is passed to both children. When the partitions $X_\text{1, leaf}$ and $X_\text{2, leaf}$ reach a leaf, we output the predictions for all the dyads in $X_\text{1, leaf} \times X_\text{2, leaf}$.
The algorithm is presented in further detail as
% \ref{alg:batch assign leaves}.
AssignLeavesOxytree
in \cref{sec:additional algorithms}.
Note that this procedure only evaluates each split rule once for each instance in the corresponding dimension, in contrast to the state-of-the-art BICTR, which performs redundant evaluations.
%For illustration purposes, consider two dyads sharing the same vertical instance, and a tree structure that happens to only have vertical splits until level 10. When predicting the two interactions, the original algorithm will evaluate the first 10 split rules twice, while we do it only once and only separate the dyads after level 10.
The resulting improvement in complexity is as follows.

Let $n_\text{test}$ and $n_\text{train}$ respectively be the numbers of test and training instances in each dimension.
The computational complexity of the Oxytrees batch inference procedure is given by
$\Theta(n_\text{test}^2)$ if $n_\text{test} \in \Omega(n_\text{train})$, and by $\Theta(n^2_\text{test} \log(n_\text{train}))$ otherwise.
%\begin{equation}
%    \Theta(n_\text{test}^2)\text{ if }n_\text{test} \in \Omega(n_\text{train})\text{, otherwise }\Theta(n^2_\text{test} \log(n_\text{train})).
%\end{equation}
%
%\begin{align}
%    \label{eq:complexity batch inference}
%    %\text{\ref{alg:batch assign leaves}}(n_\text{train},\; n_\text{test})
%    \text{\ref{alg:batch assign leaves}}
%        \in \begin{cases}
%            \Theta(n_\text{test}) & \text{if }n_\text{test} \in \Omega(n_\text{train})\text{,}\\
%            \Theta(n_\text{test} \log(n_\text{train})) & \text{otherwise.}
%        \end{cases}
%\end{align}
%
For BICTR, the complexity is $\Theta(n^2_\text{test} \log(n_\text{train}))$ in all cases.
The most common scenario is $n_\text{test} \in \Theta(n_\text{train}) \subset \Omega(n_\text{train})$ (e.g. separating a fraction of the training data for testing, or performing k-fold cross validation). Hence, the batch inference procedure is expected to run $\log (n_\text{train})$ times faster in most situations, given sufficiently large $n_\text{train}$.
These results are derived thoroughly in
% \cref{sec:complexity analysis}
appendix D2
and empirically verified in \cref{sec:empirical complexity}.


%features in the binary transformation approach.
%\subsection{Output Function of Oxytrees}
\subsection{Generalized Output Function}
\label{sec:paulownias output function}

Accurate predictions require a careful trade-off between leaf sizes that are small enough to capture local patterns and large enough to generalize well.
Leaf models are a self-adaptive middle ground in this trade-off, allowing leaves to be larger and faster to train while maintaining representation power.
Oxytrees use Regularized Least Squares with the Kronecker product kernel (RLS-Kron) as the default leaf model, a bipartite adaptation of kernel Ridge regression proposed by \citet{van_laarhoven_gaussian_2011}.
We describe below how to obtain predictions with the RLS-Kron leaf models. An illustration is provided in
% \cref{fig:oxytrees prediction}.
fig. A1.

Let $X_\text{1, leaf}$, $X_\text{2, leaf}$ and $Y_\text{leaf}$ denote the partition of the training set reaching a given leaf.
$x_\text{1, test} \in \X_1$ and $x_\text{2, test} \in \X_2$ are then test instances that were assigned to that leaf.
Given a context-specific similarity function $\phi_1: \X_1^2 \to \R$, we represent $x_\text{1, test}$ as a homonymous vector $\phi_\text{1, test}$ so that $\phi^i_\text{1, test} = \phi_1(x_\text{1, test},\; X^i_\text{1, leaf})$. 
Instances $x_\text{2, test}$ are represented analogously, and multiple individual representations are gathered in matrices $\Phi_\text{1, test}$ and $\Phi_\text{2, test}$.
%
Finally, the predicted outputs $\hat Y^{ij}$ between $\Phi^i_\text{1, test}$ and $\Phi^j_\text{2, test}$ are collectively obtained as
\begin{equation}
    \label{eq:rlskron output}
    \hat Y = \Phi_\text{2, test} \; W_\text{leaf} \;\Phi\T_\text{1, test}
    \text{.}
\end{equation}
The matrix $W_\text{leaf}$ is composed of learned linear coefficients and is determined by
\begin{equation}
    \label{eq:rlskron w}
    W_\text{leaf} = U_2
    \;[
        \Lambda
        \odot (U_2\T \; Y_\text{leaf} \; U_1)
    ]\;
    U_1\T
    \text{,}
\end{equation}
%
in which $U_1$ and $U_2$ are the matrices of eigenvectors of $\Phi_\text{1, leaf}$ and $\Phi_\text{2, leaf}$, respectively. $\Phi_\text{1, leaf}$ and $\Phi_\text{2, leaf}$, as suggested by their labels, are the symmetric similarity matrices corresponding to $X_\text{1, leaf}$ and $X_\text{2, leaf}$. The matrix $\Lambda$ has the same dimensionality as $Y_\text{leaf}$ and is defined in terms of the regularization parameter $\alpha$ and the vectors of eigenvalues $\lambda_1$ and $\lambda_2$ corresponding again to $\Phi_\text{1, leaf}$ and $\Phi_\text{2, leaf}$:
%$\Lambda^{ij} = ({\alpha + \mathbf{\lambda}_1^i \mathbf{\lambda}_2^{j}})^{-1}$.
%
\begin{equation}
    \textstyle
    %\Lambda^{ij} = \frac{1}{\alpha + \mathbf{\lambda}_1^i \mathbf{\lambda}_2^{j}}
    \Lambda^{ij} = ({\alpha + \mathbf{\lambda}_1^i \mathbf{\lambda}_2^{j}})^{-1}
    \text{.}
\end{equation}

% An interpretation for this approach is as follows. As a property of the Kronecker product, \cref{eq:rlskron output} can also be expressed as
% \begin{equation}
%     \text{vec}(\tilde Y)
%         = (\Phi_\text{1, test} \otimes \Phi_\text{2, test})
%         \; \text{vec}(W_\text{leaf})
%     \text{.}
% \end{equation}
% Therefore, RLS-Kron corresponds to learning a linear transformation on the dyad-dyad kernel matrix defined as $\Phi_\text{test} = \Phi_\text{1, test} \otimes \Phi_\text{2, test}$.
% $\Phi$ represents defining a similarity metric $\phi(d_a,\;d_b)$ between two dyads $d_a=(\x_{1a},\;\x_{2a})$ and $d_b=(\x_{1b},\;\x_{2b})$ as the product of their domain-specific similarities:
% \begin{equation}
%     \phi(d_a,\;d_b) = \phi_1(\x_{1a},\; \x_{1b}) \; \phi_2(\x_{1a},\; \x_{1b})
%     \text{.}
% \end{equation}

%An illustration of this idea is presented by \cref{fig:procedure}.
%A more complete description of RLS-Kron is available in the \cref{sec:rlskron}.



\section{Experimental Setup}
\label{sec:experimental setup}

%In this section, we describe the validation procedure and materials used in this work. First, we present the evaluation protocol utilized, followed by datasets and comparison methods. All the code and data necessary to reproduce our results is publicly available (see Data and Materials Availability section at the end of this article).

%In this section, we describe the validation procedure and materials used in this work.


%\subsection{Validation Procedure of Bipartite Learning Problems}
\subsection{Validation Procedure and Evaluation}
\label{sec:bipartite validation}

There are two ways of splitting bipartite datasets into \emph{training and test} sets to validate machine learning models: instance-wise and dyad-wise splits.
We analyze results for all split types by first performing an instance-wise split in each dimension, resulting in four partitions of the dataset, and then performing a dyad-wise split on the partition with training rows and training columns.
This results in the four test sets defined below and illustrated in \cref{fig:train test split}. L and T stand for learning (i.e., training) and test, respectively. D stands for dyads, referring to a dyad-wise split.
The naming scheme is inspired by \citet{schrynemackers_protocols_2013,pliakos_global_2018}.
%
\begin{itemize}
    % \item \textbf{Learning set (LD):}
    % $(X_\text{1, learning},\;X_\text{2, learning}, \;Y_\text{learning dyads})$
    % \item \textbf{Transductive test set (TD):}%
    % $(X_\text{1, learning},\;X_\text{2, learning}, \;Y_\text{test dyads})$
    % \item \textbf{Semi-inductive test sets:}
    % \begin{itemize}
    %     \item \textbf{Column-inductive test set (LT):}
    %         $
    %             (X_\text{1, learning},
    %             \;X_\text{2, test},
    %             \;Y_\text{learning rows, test columns})
    %         $
    %     \item \textbf{Row-inductive test set (TL):}
    %         $
    %             (X_\text{1, test},
    %             \;X_\text{2, learning},
    %             \;Y_\text{test rows, learning columns})
    %         $
    % \end{itemize}
    % \item \textbf{Inductive test set (TT):}
    % $(
    %     X_\text{1, test},
    %     \;X_\text{2, test},
    %     \;Y_\text{test rows, test columns}
    % )$
    \item \textbf{Learning set (LD):}
    %$(X_\text{1, learning},\;X_\text{2, learning}, \;Y_\text{learning dyads})$
    $(X_\text{1L},\;X_\text{2L}, \;Y_\text{LD})$
    %\item \textbf{Transductive test set (TD):} $(X_\text{1, learning},\;X_\text{2, learning}, \;Y_{\begin{subarray}{l}\text{test}\\\text{dyads}\end{subarray}})$
    \item \textbf{Transductive test set (TD):}
    $(X_\text{1L},\;X_\text{2L}, \;Y_{\text{TD}})$
    \item \textbf{Semi-inductive test sets:}
    \begin{itemize}
        \item \textbf{Column-inductive test set (LT):}
            $
                (X_\text{1L},
                \;X_\text{2T},
                \;Y_\text{LT})
            $
        \item \textbf{Row-inductive test set (TL):}
            $
                (X_\text{1T},
                \;X_\text{2L},
                \;Y_\text{TL})
            $
    \end{itemize}
    \item \textbf{Inductive test set (TT):}
    $(
        X_\text{1T},
        \;X_\text{2T},
        \;Y_\text{TT}
    )$
\end{itemize}

\begin{figure}[tbh]
  %\begin{minipage}[c]{0.5\textwidth}
    \def\svgscale{.3}
    \centering
    %\input{figures/illustrations/split.pdf_tex}
    \input{figures/illustrations/oxytrees_figures_train_test_v2.pdf_tex}
  %\end{minipage}\hfill
  %\begin{minipage}[l]{0.5\textwidth}
    \caption{Diagram of the test (T) and training (i.e. learning, L) sets in the bipartite context (see \cref{sec:bipartite validation}).}
    \label{fig:train test split}
  %\end{minipage}
\end{figure}
%Dyad-wise splits intrinsically imply a \emph{transductive}~\cite{chapelle_semi-supervised_2006,lu_link_2011,zhou_semi-supervised_2021} learning setting.
%%
%Therefore, although common in the literature, dyad-wise splitting fails to evaluate the ability of a model to generalize to unseen instances (i.e. inductiveness), as pointed out by \citet{park_flaws_2012,schrynemackers_protocols_2013,pahikkala_toward_2015}.
%Various names for this partitioning procedure and test sets have been proposed in the literature~\cite{park_flaws_2012,schrynemackers_protocols_2013,pahikkala_toward_2015,liu_neighborhood_2016,pliakos_global_2018}. We opt for the concise two-letter acronyms above.

For \emph{cross-validation}, each of the two domains was split into $k$ folds, using label stratification to encourage consistent fold densities. Each (row fold)-(column fold) pair can represent a TT set, yielding $k^2$ possible TT sets in total. The corresponding semi-inductive test sets were also collected for each TT set, and the training set and transductive test set were then generated.
The value of $k$ was adjusted for each learning task, using more folds for smaller datasets (\cref{tab:datasets}).
%\Cref{fig:cv} illustrates this procedure.


A percentage of the positive annotations was masked (replacing ones with zeros) to generate $Y_\text{TD}$.
We refer to this percentage as the \emph{positives masking percentage} (PMP).
All the missing annotations in the learning set (zeros that were not generated by masking) were also used as negative annotations for evaluation in the transductive setting. We explored four values for PMP: $0\%$, $25\%$, $50\%$, and $75\%$.

% \subsection{Evaluation Measures and Statistical Tests}
% \label{sec:evaluation}

Regarding \emph{evaluation measures}, we have employed the area under the precision and recall curve (AUPRC) and the area under the ROC curve (AUROC), following  \citet{schrynemackers_protocols_2013,pliakos_global_2018,pliakos_drug-target_2020,liu_drug-target_2022}. 
%
Moreover, we have employed the Friedman test, followed by the post-hoc Nemenyi test, as suggested by \citet{demsar_statistical_2006}.
%Graphically, the results are displayed in a critical difference diagram where methods placed in lower ranks are associated to better performance. Further, methods connected by a horizontal bar are not statistically significantly different.
%
We report results for the semi-inductive settings altogether, averaging ranks of corresponding TL and LT scores.


\subsection{Datasets}
\label{sec:datasets}

\Cref{tab:datasets} describes the datasets used in this work.
As can be seen, we have employed a total of 15 datasets from several biological domains.
All features correspond to similarities measured between instances of the same input space.
%For example, the features of dataset DPI-E associated to $X_1$ (drugs) correspond to the similarities among instances from $X_1$ (drugs). Its counterpart, $X_2$, contains similarities among the instances from $X_2$, enzymes.
Feature matrices in the test sets represent similarities with the training instances (avoiding data leakage).

\begin{table}[ht]
\centering
\small
\setlength{\tabcolsep}{1mm}  % allowed by author kit
\begin{tabular}{llllrrl}
%\begin{tabular}{lllll}
\toprule
Name & Domain & Size & Folds & Density \\
%Name & Domain & Dimensionality & Folds & Density \\
\midrule
DPI-N & Drug-NR & $26 \times 54$ & $4 \times 4$ & 6.4\% \\
TE-piRNA & TE-piRNA & $60 \times 93$ & $4 \times 4$ & 3.9\% \\
lnc-mi & LncRNA-miRNA & $44 \times 218$ & $3 \times 3$ & 10\% \\
DPI-G & Drug-GPCR & $95 \times 223$ & $3 \times 3$ & 3.0\% \\
Davis & Inhibitor-Kinase & $68 \times 442$ & $3 \times 3$ & 5.0\% \\
lnc-D & LncRNA-Disease & $156 \times 241$ & $3 \times 3$ & 6.6\% \\
lnc2Cancer & LncRNA-Cancer & $367 \times 106$ & $3 \times 3$ & 7.3\% \\
DPI-I & Drug-Ion channel & $204 \times 210$ & $2 \times 2$ & 3.5\% \\
MiRNA-D & MiRNA-Disease & $462 \times 252$ & $2 \times 2$ & 11\% \\
ERN & Gene-TF & $1164 \times 154$ & $2 \times 2$ & 1.8\% \\
SRN & Gene-TF & $1821 \times 113$ & $2 \times 2$ & 1.8\% \\
NPInter & LncRNA-Protein & $586 \times 446$ & $2 \times 2$ & 18\% \\
DPI-E & Drug-Enzyme & $664 \times 445$ & $2 \times 2$ & 1.0\% \\
KIBA & Inhibitor-Kinase & $2111 \times 229$ & $2 \times 2$ & 20\% \\
miRTarBase & MiRNA-Gene & $1873 \times 415$ & $2 \times 2$ & 7.1\% \\
\bottomrule
\end{tabular}
\caption{Summary of the datasets used in this study. Density refers to the percentage of positive annotations in each dataset. NR = nuclear receptor; TE = transposable element; TF = transcription factor.}
\label{tab:datasets}
\end{table}

DPI-N, DPI-G, DPI-I, DPI-E, SRN, ERN, Davis and KIBA are widely employed in the literature \cite{schrynemackers_classifying_2015,pliakos_network_2019,huang_moltrans_2020}. We have preprocessed NPInter and mirTarBase in our preliminary work~\cite{ilidio_fast_2024}. We have also collected five more datasets, TE-Pirna, LncRNA-miRNA, lncRNA-disease, lncRNA-cancer, and miRNA-disease, which, to the best of our knowledge, have never been used as benchmarks for interaction prediction in general.
%
Lnc2Cancer contains lncRNA-cancer interactions \cite{gao2021lnc2cancer}\footnote{http://bio-bigdata.hrbmu.edu.cn/lnc2cancer/}.
%
The datasets LncRNA-disease, LncRNA-miRNA and miRNA-disease were made available in \cite{sheng2023data}, where the authors collected data from several repositories.
%
TE-piRNA contains interaction between transposable elements (TEs) and piwi-RNAs (piRNAs)~\cite{dos2024transposable}.


\subsection{Comparison Methods}
\label{sec:baselines}


We present in \cref{tab:method-summary} a list of the comparison methods used in this work. We selected prominent methods from the literature that i) were able to infer interactions for new instances, and ii) were applicable across different bipartite learning problems. We provide more details about them in
% \cref{sec:baselines details}.
appendix E.
As previously mentioned, the most common deep learning approaches do not meet criterion (ii). We thus include a baseline multi-layer perceptron (MLP)~\cite{hinton1990connectionist} under the data-based global adaptation~\cite{schrynemackers_protocols_2013} to still represent this model category. Following previous literature~\cite{huang_moltrans_2020,dehghan_tripletmultidti_2023}, we performed random undersampling of negative annotations, resulting in balanced training datasets for the MLP baseline.

\begin{table}[tbh]
\centering
\small
\setlength{\tabcolsep}{1mm}
\begin{tabular}{lll|cccc}
\toprule
\textbf{Method}
    & \textbf{Base technique}
    & \textbf{Adapt.}
    & \textbf{TF}
    & \textbf{Ex}
    & \textbf{EB}
    & \textbf{In}
    % & \makecell[l]{\textbf{Adapt.} \\ \textbf{strategy}}
    % & \rotatebox[origin=l]{90}{\textbf{Tuning-free}}
    % & \rotatebox[origin=l]{90}{\textbf{Interpretable}}
    % & \rotatebox[origin=l]{90}{\textbf{Estimator-based}}
    % & \rotatebox[origin=l]{90}{\textbf{Inductive}}
    \\
\midrule
\textbf{MLP}
    & Deep learning
    & Global
    &
    &
    &
    & \checkmark
    \\
\textbf{RLS-avg}
    & Linear regression (LR)
    & Local
    &
    & \checkmark
    &
    &
    \\
\textbf{RLS-Kron}
    & LR + Kronecker product
    & ---
    &
    & \checkmark
    & \checkmark
    & \checkmark
    \\
\textbf{BLMNII}
    & LR + Nearest Neighbors
    & Local
    &
    & \checkmark
    &
    & \checkmark
    \\
\textbf{NRLMF}
    & Matrix factorization
    & ---
    &
    &
    & \checkmark
    & \checkmark
    \\
\textbf{WkNNIR}
    & Nearest-Neighbors
    & Local
    & \checkmark
    &
    & 
    & \checkmark
    \\
\textbf{BICTR}
    & Decision trees
    & ---
    &
    & \checkmark
    & \checkmark
    & \checkmark
    \\
\midrule
\textbf{Oxytrees}%[Deep, YR]}
    & Decision trees
    & ---
    & \checkmark
    & \checkmark
    & \checkmark
    & \checkmark
    \\
% \textbf{Oxytrees[Logistic]}
%     & Tree-based
%     & \checkmark
%     & \checkmark
%     & \checkmark
%     & \checkmark
%     \\
% \textbf{Oxytrees[RLS-Kron]}
%     & Tree-based
%     & \checkmark
%     & \checkmark
%     & \checkmark
%     & \checkmark
%     \\
\bottomrule
\end{tabular}
\caption{%
    Summary of methods evaluated in our experiments. References:
    MLP~\cite{hinton1990connectionist}; 
    RLS-\{avg, Kron\}~\cite{van_laarhoven_gaussian_2011};
    BLMNII~\cite{mei2013drug};
    NRLMF~\cite{liu_neighborhood_2016};
    WkNNIR~\cite{liu_drug-target_2022};
    BICTR~\cite{pliakos_drug-target_2020}.
    See
    % \cref{sec:baselines details}
    appendix E
    for hyperparameters of each method.
    Legend: Adapt. = bipartite adaptation strategy; TF = hyperparameter tuning-free; Ex = explainable; EB = Estimator-based adaptation (as opposed to data-based); In = inductive.
}
\label{tab:method-summary}
\end{table}



Our proposal does not enforce any particular ensembling technique.
Following BICTR~\cite{pliakos_drug-target_2020}, we used extremely randomized~\cite{geurts_extremely_2006} Oxytrees for this study. 

Regarding hyperparameter tuning, a nested 2 by 2 cross-validation procedure was utilized, using AUPRC to select the best hyperparameter combination for each outer fold.
The specific hyperparameter sets considered for each method are also presented in
% \cref{sec:baselines details}.
appendix E.


\section{Results and Discussion}
\label{sec:results}

\subsection{Oxytrees Against Previous Proposals}
\label{sec:literature comparison}

Inductive AUPRC and AUROC results for 0\% PMP are presented in \cref{fig:literature inductive subset}. Similar results were found for other PMP values
% (\cref{fig:literature inductive}).
(fig. F1).
Average values for AUROC and AUPRC in each dataset are presented by
% \cref{tab:auroc} and \cref{tab:auprc},
table F2 and table F1,
respectively.
%
\begin{figure*}[tbh]
    \includegraphics[width=.45\textwidth]{results/literature_methods/TT/all_datasets/critical_difference_diagrams/AUPRC (Inductive).pdf}
    \hfill
    \includegraphics[width=.45\textwidth]{results/literature_methods/TT/all_datasets/critical_difference_diagrams/AUROC (Inductive).pdf}
    \caption{Comparisons of AUPRC and AUROC of the proposed Oxytrees against previous methods for bipartite learning. The inductive setting is analysed, for 0\% positives masking percent (PMP) (see \cref{sec:bipartite validation}).
    Further results in
    % \cref{fig:literature inductive}.
    fig.~F1.
    }
    \label{fig:literature inductive subset}
\end{figure*}
%
Oxytrees had the best average ranking for all but one of the inductive settings analyzed, and no statistically significant differences were found comparing Oxytrees to the state-of-the-art.
%
We also see that Oxytrees perform consistently better than RLS-Kron alone, indicating that the tree structure provides considerable improvement.
%It suggests that RLS-Kron is better suited to model the smaller groups of instances in each leaf, rather than the task in its entirety. This is similar in essence to how a linear function can often locally approximate another function with more complicated behavior.
%
% \Cref{fig:literature semi-inductive}
Figure~F2
presents results for the semi-inductive setting. Oxytrees, BICTR and RLS-Kron were the top performers, and the differences between their scores were not statistically significant in any case.
%
For the transductive setting
% (\cref{fig:literature transductive}),
(fig.~F3),
NRLMF, BICTR and RLS-Kron had the best overall scores.
%
%These results highlight the potential of Oxytrees for bipartite learning: Oxytrees consistently displayed the best performance despite often running an order of magnitude faster than BICTR, its strongest competitor in terms of AUROC.

% 1 - Difference in performance between models at the leaf: KronRLS vs Logistic
%Comparing the two different leaf models in all settings, we observe consistent superiority of Oxytrees[RLS-Kron] against Oxytrees[Logistic]. A crucial difference between the two models is that RLS-Kron uses all combinations $\phi_1^i\phi_2^j$ between elements of the feature vectors $\phi_1$ and $\phi_2$ in each domain, while the logistic regressor uses $\phi_1$ concatenated with $\phi_2$. We argue that the superior performance of RLS-Kron is due to this extra complexity: the larger number of input features seems to still be necessary to capture the behavior of our data.

% 2 - Difference between NLRMF + Oxytrees (Deep) and BICTR performance wise
%Comparing Oxytrees[Deep, YR] against BICTR, we notice that the Oxytrees perform slightly better in almost all scenarios, and remarkably close otherwise. This is expected, since the only fundamental difference between the models is their impurity function (\cref{sec:oxytrees impurity}). However, the expressive computational advantage of Oxytrees (\cref{sec:oxytrees training} and \cref{sec:empirical complexity}) indicates Oxytrees as the best option even for scenarios where BICTR was dominant.

%  6- Discuss numbers from the tables. Say that despite our method achieving better results and expanding the state-of-the-art, this problem is still very challenging since the numbers are overall still low
% Considering the values for AUPRC and AUROC for each dataset (Tab. \ref{tab:auprc} and \ref{tab:auroc}), we notice very different results for each task, suggesting that the choice of predictive algorithm to use must take into consideration the specific scenario of the learning problem.
The scores in the inductive
% (Tab. \ref{tab:auprc} and \ref{tab:auroc})
(tables F1 and F2)
and semi-inductive
% (Tab. \ref{tab:auprc lt}, \ref{tab:auroc lt}, \ref{tab:auprc tl} and \ref{tab:auroc tl})
(tables F5-F8)
settings are considerably lower compared to the transductive setting
% (Tab. \ref{tab:auprc transductive} and \ref{tab:auroc transductive}),
(tables F3 and F4),
which highlights the general difficulty of considering new instances.


 
% 3 - Short sentence about performance of the baselines. Local vs Global. Kron is still quite strong bla bla, local methods are not doing that great. NRLMF by itself is not great either, mb that has an effect on the reconstruction...... 
% Regarding the baselines other than BICTR, RLS-avg and RLS-Kron showed the best overall performance (\cref{fig:literature inductive}), followed by RLS-avg. Their performance was only worse than Oxytrees[RLS-Kron] for AUPRC, and worse than all tree-based models for AUROC.
On the other hand, NRLMF was one of the worst performers in the inductive setting compared to the other models. This shows that the original promising results of NRLMF in the transductive setting (found by \citet{liu_neighborhood_2016,liu_lpi-nrlmf_2017,liu_predicting_2020} and reproduced by us in
% \cref{fig:literature transductive})
fig. F3).
do not translate into the inductive setting. This motivates the previously proposed combination of NRLMF and biclustering forests~\cite{pliakos_drug-target_2020}, restricting NRLMF to the transductive step of matrix completion and delegating the inductive modeling to the forest.
Furthermore, WkNNIR was not able to outperform BICTR in our experiments with 15 datasets, contrary to what was observed for drug-target interaction prediction~\cite{liu_drug-target_2022}.

% 4 - Differences between the different percentages of masking 
To analyze the robustness of each model to missing positives, we compare their scores in different PMP settings relative to the setting with PMP = 0\%
% (\cref{fig:pmp drop}).
(fig. F4).
As expected, the performances of all models decrease when we increase the PMP, with very few exceptions that we attribute to the intrinsic noise of these values.
BLMNII and RLS-avg were the models with the smaller performance drop in general, suggesting that the simplicity of local linear models makes them robust to missing labels.
The results for MLP were inconsistent, highlighting its sensitivity to the learning setting and hyperparameter choices.



\subsection{Ablation Study}
\label{sec:ablation}

We compare ensembles of Oxytrees with and without the minimum leaf size constraint and different leaf models, as indicated in \cref{tab:ablation}. We also test the interaction matrix reconstruction (YR) technique proposed by \citet{pliakos_drug-target_2020}.
% To compare across datasets, we divided the scores of each fold by the scores of the main model.
The results are presented in \cref{fig:ablation}.

Removing the leaf model (and using the average label as usual) significantly degraded the performance, even when using YR.
Using YR while enforcing minimum leaf dimensions of 5 by 5 was the worst overall strategy.
%
\begin{table}[htb]
\centering
\small
\setlength{\tabcolsep}{1mm}  % allowed by author kit

\begin{tabular}{lccc}
\toprule
% requires hyperparameter tuning, explainable, estimator-based, inductive
\textbf{Model}
    & \textbf{Leaf model}
    % & \textbf{Minimum leaf size 5 by 5}
    & \textbf{Min. 5 by 5 leaves}
    %& \textbf{$Y$ reconstruction with NRLMF}
    & \textbf{$Y$ rec.}
    \\
\midrule
\textbf{Oxytrees}
    & RLS-Kron
    & \checkmark
    &
    \\
\midrule
\textbf{[Mean]}
    & ---
    & \checkmark
    &
    \\
\textbf{[Logistic]}
    & Logistic reg.
    & \checkmark
    &
    \\
\textbf{[Deep]}
    & ---
    &
    &
    \\
\textbf{[YR]}
    & RLS-Kron
    & \checkmark
    & \checkmark
    \\
\textbf{[Mean, YR]}
    & ---
    & \checkmark
    & \checkmark
    \\
\textbf{[Deep, YR]}
    & ---
    &
    & \checkmark
    \\
\bottomrule
\end{tabular}
\caption{%
    Summary of method modifications compared in the ablation study.
    Interaction matrix reconstruction ($Y$ rec.) was performed with NRLMF~\cite{liu_neighborhood_2016}, following ~\citet{pliakos_drug-target_2020}.
    For the ``Deep" models, the trees were grown until all dyads in a node share the same label. For the ``Mean" models, the mean label was used for the prediction, as usual for decision trees.
}
\label{tab:ablation}
\end{table}

\begin{figure}
    %\includegraphics[width=0.49\linewidth]{figures/results/ablation_relative/Inductive_AUROC_00.pdf}
    \includegraphics[width=0.49\linewidth]{figures/results/ablation_relative/Inductive_AUPRC_00.pdf}
    \hfill
    \includegraphics[width=0.49\linewidth]{figures/results/ablation_relative/Inductive_AUPRC_50.pdf}
    \caption{
    Performance comparison of ensembles of Oxytrees using different components (\cref{tab:ablation}). To compare across datasets, the scores are divided by the score of the main model (Oxytrees). Each point represents a cross-validation fold, and the mean value is presented in the white boxes.
    Asterisks indicate significance in comparison to Oxytrees ($p<0.05$, Wilcoxon signed-rank), averaging folds for each dataset.
    Remaining results are presented in
    % \cref{fig:ablation extra}.
    fig. F6.
    }
    \label{fig:ablation}
\end{figure}


%\subsection{Empirical Computational Complexity Analysis}
\subsection{Empirical Complexity Analysis}
\label{sec:empirical complexity}
%\subsection{Run time comparison}

In this analysis, we compared the training and inference times of Oxytrees and BICTR, for different data sizes.
For each size $n_1$, we generated three $n_1$ by $n_1$ matrices with pseudo-random values to serve as $X_1$, $X_2$, and $Y$. We then built a single tree of the BICTR and Oxytrees ensembles for each $n_1$.
We also built an Oxytree until each leaf contained a single instance (Oxytree[Deep], \cref{tab:ablation}), to isolate the effect of the larger leaves.
To measure inference times, we built a completely random tree for each $n_1$ and used it twice to assign a leaf to each instance: through the BICTR procedure, and through the Oxytrees procedure (\cref{sec:batch inference}).
The results are shown in \cref{fig:empirical complexity}. The estimated complexities follow the theoretical expectations of \cref{sec:oxytrees training} (with $m = n_1$) and \cref{sec:batch inference}.
%: training in $\Theta(n^3)$ for Oxytrees and $\Theta(n^3 \log (n))$ for BICTR; inference in $\Theta(n^2)$ for Oxytrees and $\Theta(n^2\log(n))$ for BICTR.
Plotting times of Oxytrees against BICTR
% (\cref{fig:complexity extra}),
(fig. F5),
we measured that Oxytrees trained $35.49 \pm 0.75$ times faster
%($p < 4 \cdot 10^{-8}$, Wald)
and predicted $9.98 \pm 0.84$ times faster ($p < 4 \cdot 10^{-5}$, Wald).

% Results Oxytrees vs. BICTR
% fit_time: 35.49 \pm 0.75 ($p = 3.924593974264197e-08$)
% predict_time: 9.98 \pm 0.84 ($p = 3.722418598428486e-05$)


%For all cases, we generate artificial training sets $X_1$, $X_2$ and $Y$ by filling three $n \times n$ matrices with uniform pseudo-random values between 0 and 1. \Cref{fig:fit complexity} compares extremely randomized versions of a single Oxytree[RLS-Kron], a single Oxytree[Deep] and a single PBCT[Deep] for multiple values of $n$. The [Deep] suffix means that the trees were grown to their maximum depth, whereas Oxytree[RLS-Kron] used a minimum leaf-size of 5 by 5 as a stopping criterion.
%\Cref{fig:predict complexity} compares the inference times of the proposed \cref{alg:batch assign leaves} used by Oxytrees with the previous individual inference used by PBCT~\cite{pliakos_global_2018}. For each $n$ in \cref{fig:predict complexity}, a completely random tree is used, and the same tree is used for both inference methods.
%
\begin{figure}[tbh]
    \centering
    %\begin{subfigure}{.47\textwidth}
    %    \includegraphics[width=\linewidth]{figures/results/empirical_complexity/training_empirical_complexity_2.pdf}
    %    \caption{Training time as a function of the number of instances.}
    %    \label{fig:fit complexity}
    %\end{subfigure}
    %\hfill
    %\begin{subfigure}{.47\textwidth}
    %    \includegraphics[width=\linewidth]{figures/results/empirical_complexity/predicting_empirical_complexity_2.pdf}
    %    \caption{Inference time as a function of the number of instances.}
    %    \label{fig:predict complexity}
    %\end{subfigure}
    
    \includegraphics[width=0.49\linewidth]{figures/results/empirical_complexity_plots/empirical_complexity_fit.pdf}
    \includegraphics[width=0.49\linewidth]{figures/results/empirical_complexity_plots/empirical_complexity_leaf_assign.pdf}
    % \includegraphics{figures/results/empirical_complexity_plots/empirical_complexity_fit.pdf}
    % \includegraphics{figures/results/empirical_complexity_plots/empirical_complexity_leaf_assign.pdf}
    
    \caption{Empirical complexity analysis. Single trees were applied to artificial datasets of different dimensions (\cref{sec:empirical complexity}).
    We used the last 10\% of the points of each curve to approximate the asymptotic complexity as $\Theta (n^\alpha)$. $\alpha$ ($\pm$ standard dev.) is estimated as
    %the slope in the log-log space.
    the slope of the linear regression in the log-log space.
    %BICTR leaf assignment terminated earlier due to memory limitations.
    Further results in
    % \cref{fig:complexity extra}.
    fig. F5.
    }
    \label{fig:empirical complexity}
\end{figure}

%The empirical results follow the theoretical expectations: Oxytrees have training time of order $\Theta(n^3)$ (as stated by \cref{eq:complexity oxytree} with $m=n$) and inference time of order $\Theta(n^2)$ (as stated by \cref{eq:complexity batch inference}). The order of complexity of PBCT training is asymptotically between $n^3$ and $n^4$, as suggested by \cref{eq:complexity pbct}, while its inference complexity is between $n^2$ and $n^3$, which is compatible with the $n^2\log \, n$ complexity discussed in \cref{sec:batch inference} and \cref{sec:predict complexity analysis}.


\subsection{Dependency on Leaf Size}
\label{sec:leaf size}

We assess the effect of setting minimum leaf dimensions on the predictive performance of biclustering forests, comparing Oxytrees with RLS-Kron leaf models, Oxytrees without leaf models and BICTR. The results for the inductive setting are shown in \cref{fig:leaf size}.
Similar results were found for the other settings
% (\cref{sec:additional results}).
(appendix F).
The performance of forests without leaf models drops when the leaves are enforced to be larger, especially for leaf dimensions larger than 20 by 20. On the other hand, Oxytrees with leaf models keep their performance stable even with larger leaves, with a slight score increase up until 20 by 20 leaves.
Therefore, adding RLS-Kron as leaf models enables shallower trees.
%while improving performance by allowing more complex relationships to be modeled in each leaf.

\begin{figure}[tbh]
    \centering
    % \includegraphics[width=\linewidth]{figures/results/leaf_size/Inductive_AUPRC_ 0.pdf}
    % \includegraphics[width=\linewidth]{figures/results/leaf_size/Inductive_AUROC_ 0.pdf}
    \includegraphics{figures/results/leaf_size/Inductive_AUPRC_ 0.pdf}
    \includegraphics{figures/results/leaf_size/Inductive_AUROC_ 0.pdf}
    \caption{
    Performance of biclustering forests as a function of minimum leaf dimensions, relative to the initial performance with 2 by 2 leaves.
    %of BICTR, Oxytrees with RLS-Kron leaf models, and Oxytrees without leaf models ()
    Values in the $x$ axis represent the minimum number of instances in each dimension.
    %(e.g. 20 means leaves equal or larger than 20 by 20, totalling at least 400 samples per leaf).
    Markers indicate mean and standard deviation over 15 datasets.
    }
    \label{fig:leaf size}
\end{figure}


\subsection{Dependency on the Number of Trees}
\label{sec:number of trees}

We perform a bootstrapping procedure to evaluate how the performance of each biclustering forest varies with the number of trees.
For each forest, we started by building a total of 200 trees.
Then, we progressively added the outputs of each tree, calculating the resulting score in each step. We repeated the summing 50 times, with a random tree order each time.
From the plot of scores vs. number of trees, we used linear interpolation to estimate the number of trees required to achieve 98\% of the final score of 200 trees. The measurements for PMP 0\% are reported in \cref{fig:number of trees}.
% The large number of points below the $x=y$ line demonstrates that ensembles of Oxytrees require significantly less trees than BICTR to approximate their full performance potential, as confirmed by a Wilcoxon signed-rank test.
Oxytrees required significantly fewer trees ($42\%\pm 3\%$) relative to BICTR ($p < 2 \cdot 10^{-4}$, Wilcoxon).
%($42\pm 7$)\%
% We report the score of each step relative to the final score with all 200 trees (\cref{fig:number of trees}).

\begin{figure}[tb]%!]
    % \includegraphics[width=.48\linewidth]{figures/results/number_of_trees/Inductive_AUROC_0_BICTR_Oxytrees.pdf}
    % \hfill
    % \includegraphics[width=.495\linewidth]{figures/results/number_of_trees/Inductive_AUPRC_0_BICTR_Oxytrees.pdf}
    \includegraphics{figures/results/number_of_trees/Inductive_AUROC_0_BICTR_Oxytrees.pdf}
    \includegraphics{figures/results/number_of_trees/Inductive_AUPRC_0_BICTR_Oxytrees.pdf}
    
    \caption{Comparison between Oxytrees and BICTR in terms of the expected number of trees to achieve 98\% of the score of 200 trees. Each fold is plotted, with different styles for each dataset.
%    The p-value for the comparison
%    (Wilcoxon signed-rank)
%   is displayed in the legend, based on the averaged scores per dataset.
     Remaining results can be found in
     % \cref{fig:number of trees extra}.
     fig. F8.
    }
    \label{fig:number of trees}
\end{figure}



\section{Conclusion and Future Work}
\label{sec:conclusion}

We proposed an ensemble of proxy-based biclustering model trees: Oxytrees.
%Oxytrees use lower-dimensional proxies of rows and columns of the interaction matrix as well as linear models in each leaf to significantly speed up training. In addition, we introduce a faster leaf assignment algorithm to also speed up prediction.
%
%employs a novel tree building and batch inference algorithms, as well as constructs a model at each leaf to make its predictions.
%
%Using 15 datasets and different percentages of positive annotations masked, we compared our method against several prominent baselines.
% including BICTR~\cite{pliakos_drug-target_2020}, the current state-of-the-art for inductive drug-target interaction prediction.
%Oxytrees had the best average rank among state-of-the-art techniques for inductive bipartite learning.
Oxytrees had competitive results among state-of-the-art techniques for inductive bipartite learning.
Against the state-of-the-art biclustering forest, BICTR, Oxytrees were $35.49 \pm 0.75$ times faster to train ($p < 4 \cdot 10^{-8}$) and predicted $9.98 \pm 0.84$ times faster ($p < 4 \cdot 10^{-5}$). They also required up to 50\% less trees to achieve 98\% of the performance of 200 trees ($p < 2 \cdot 10^{-4}$).
%Our results showed that Oxytrees had superior or competitive predictive performance in all settings, besides having considerably lower computational complexity compared to previous biclustering trees.
%
%Regularized linear regression with the Kronecker product kernel (RLS-Kron) was the best leaf model according to our results, compared to global single output logistic regression.
The model comparisons were mostly consistent across different masking percentages of positive annotations.
%
We 
%introduce a new Python package,
%texttt{bipartite\_learn},
%texttt{},
provide datasets and efficient implementations of the algorithms compared in this study, in an accessible API based on Scikit-Learn.
%
%
As future work, we would like to explore the potential of Oxytrees for regressive bipartite learning, such as drug-target affinity prediction, and compare our results with deep learning techniques that work on the raw unstructured data (such as protein sequences and molecular structures).

\FloatBarrier

\iffalse

\bmhead{Data and Materials Availability}
\label{sec:availability}
All the data and scripts needed to reproduce our results are made available in a supplementary git repository\footnote{\url{https://github.com/pedroilidio/oxytrees2025}}.  
%
All the compared algorithms are implemented and made available through our Python library \texttt{bipartite\_learn}\footnote{\url{https://github.com/pedroilidio/bipartite_learn}}.  

\bmhead{Acknowledgements}

\fi
%\FloatBarrier

\section*{Acknowledgements}
The work received funding from the Flemish Government (AI Research Program), VOEWIAI02.
The authors would like to thank FWO, 11A7U26N, 1235924N and 1S38025N.
R.C. thanks the Brazilian National Council for Scientific and Technological Development (CNPq), grant number 300934/2025-4, and the São Paulo Research Foundation (FAPESP), Brasil, Process Number \#2022/02981-8.

\bibliography{bibliography}
\vfill

\iftrue
\pagebreak
\begin{appendices}

\crefalias{section}{appendix}
\crefalias{enumi}{definition}
\counterwithin{figure}{section}
\counterwithin{table}{section}
\renewcommand\thefigure{\thesection\arabic{figure}}
\renewcommand\thetable{\thesection\arabic{table}}

\renewcommand\thetheorem{\thesection\arabic{theorem}}

\section{Additional Illustrations}
\label{sec:additional illustrations}
\FloatBarrier

%\begin{figure}[tbh]
%    \centering
%    \def\svgscale{.3}
%    \input{figures/illustrations/lmo.pdf_tex}
%    \caption{Caption}
%    \label{fig:lmo}
%\end{figure}

\begin{figure}[tbhp]
    \centering
    \def\svgscale{.3}
    \input{figures/illustrations/oxytree_prediction.pdf_tex}
    \caption{Illustration of the inference procedure of an Oxytree with the RLS-Kron as the leaf model. The $\phi$ feature vectors contain similarities of the test instance with the training instances. We use $\otimes$ to denote the Kronecker product and $\odot$ to denote the Hadamard (element-wise) product. $W_\text{leaf}$ represent the matrix of linear coefficients learned during training for the leaf node illustrated.}
    \label{fig:oxytrees prediction}
\end{figure}


\FloatBarrier

\section{Additional Algorithms}
\label{sec:additional algorithms}
\FloatBarrier

\begin{function}[tb]
    \caption{FindBestSplitOxytree($X_1$, $X_2$, $Y$): The split search procedure of Oxytrees.}
    \label{alg:find best split bgso}
    
    \KwIn{The partition of the training set that reached a given tree node}
    \KwOut{The best split rule among the evaluated split candidates}
    \BlankLine
    Determine the set $S_1$ of horizontal split rules to consider\;
    
    Determine the set $S_2$ of vertical split rules to consider\;
    \BlankLine
    
    Using $Y$, build the proxy matrices $\tilde Y_1$ and $\tilde Y_2$ (\cref{eq:proxies})\;
    %\tcp{$\tilde Y_1^i$ summarizes all the interactions of $X_1^i$.}
    %\tcp{$\tilde Y_2^j$ summarizes all the interactions of $X_2^j$.}
    %\tcp{For our specific impurity in \cref{eq:i bgso}, use \cref{eq:proxies}.}
    \BlankLine

    %$s^\star \gets null$  \tcp*{Initialize the best split}
    $s^\star \gets null$  \tcp*{The best split}
    $\Delta I^\star \gets 0$  \tcp*{The best $I$ improvement}
    \BlankLine
    
    \ForEach {horizontal split rule $s_1 \in S_1$}{
        Determine $\Delta I (s_1,\; X_1,\; \tilde Y_1)$\label{ln:y1}\;
        \BlankLine
        \uIf{$\Delta I (s_1, X_1, \tilde Y_1) > \Delta I^\star$}{
            $s^\star \gets s_1$ 
        }
    }
    \ForEach {vertical split rule $s_2 \in S_2$}{
        Determine $\Delta I (s_2,\; X_2,\; \tilde Y_2)$\label{ln:y2}\;
        \BlankLine
        \uIf{$\Delta I (s_2, X_2, \tilde Y_2) > \Delta I^\star$}{
            $s^\star \gets s_2$ 
        }
    }
    \BlankLine
    \Return $s^\star$
\end{function}



% \begin{function}[tb]
%     \caption{BuildTree($X$, $Y$): General procedure for building a decision tree}
%     \label{alg:dt build}
%     
%     \KwIn{The training data}
%     \KwOut{A tree node pointing to its descendants}
%     \SetKwFunction{rightchild}{RightChild}
%     \SetKwFunction{leftchild}{LeftChild}
%     \SetKwFunction{node}{Node}
%     \SetKwFunction{leaf}{Leaf}
%     \SetKwFunction{self}{BuildTree}
%     \SetKwFunction{findbestsplit}{FindBestSplit}
%     \BlankLine
% 
%     \If{Stopping criteria are met}{
%         Create a leaf node \leaf\;
%         \Return{\leaf}
%     }
%     
%     \BlankLine
%     $s^\star \gets$ \findbestsplit($X$, $Y$)
%     \tcp*{Determine the split rule to be used}
%     \BlankLine
%     
%     Create a new node \node containing $s^\star$\;
%     
%     \BlankLine
%     Use $s^\star$ to obtain dataset partitions ($X_\text{left}$, $Y_\text{left}$) and ($X_\text{right}$, $Y_\text{right}$)\;
%     
%     \BlankLine
%     \leftchild $\gets$ \self($X_\text{left}$, $Y_\text{left}$)
%     \tcp*{Create children nodes}
%     \rightchild $\gets$ \self($X_\text{right}$, $Y_\text{right}$)\;
%     
%     \BlankLine
%     \Return{\node}
% \end{function}

\begin{function}[tb]
    \caption{FindBestSplit($X$, $Y$): A generic split search procedure for greedy decision trees}
    \label{alg:find best split dt}
    
    \KwIn{The partition of the training set that reached a given tree node}
    \KwOut{The best split rule among the evaluated split candidates}
    \BlankLine
    %\SetKwFunction{SplitQuality}{SplitQuality}
    %\SetKwFunction{BestSplit}{BestSplit}
    %\SetKwFunction{SplitDataset}{SplitDataset}
    %\tcp{Index of the feature ($X$ column) used in the best split found}
    %Initialize \BestSplit object\;
    %Calculate $\savg{Y^{i}}_i$ and $\savg{(Y^{i})^2}_i$\;
    %$S_l \gets 0$ \tcp*{The same but for $Y_l$}
    % $f^* \gets 0$ \tcp*{best split's feature}
    % $\Delta I^* \gets 0$ \tcp*{best split's quality}
    % $t^* \gets 0$ \tcp*{best split's threshold}
    Determine the set $S$ of split rules to consider.
    \BlankLine

    $s^\star \gets \emptyset$  \tcp*{Initialize the best split}
    \BlankLine
    
    \ForEach {split rule $s \in S$}{
        Determine $\Delta I (s, X, Y)$\;
        \BlankLine
        \If{$\Delta I (s, X, Y) > \Delta I (s^\star, X, Y)$}{
            $s^\star \gets s$ 
        }
    }
    \BlankLine
    \Return $s^\star$
\end{function}

\begin{function}[tb]
    \caption{FindBestSplitBICTR($X_1$, $X_2$, $Y$): Split search procedure used by methods based on PBCT~\cite{pliakos_global_2018}, including the main competitor BICTR~\cite{pliakos_drug-target_2020}.}
    \label{alg:find best split gmo}
    
    \KwIn{The partition of the training set that reached a given tree node}
    \KwOut{The best split rule among the evaluated split candidates}
    \BlankLine
    Determine the set $S_1$ of horizontal split rules to consider.\;
    
    Determine the set $S_2$ of vertical split rules to consider.\;
    \BlankLine
    
    $s^\star \gets null$  \tcp*{Initialize the best split}
    $\Delta I^\star \gets 0$  \tcp*{Initialize the best impurity improvement}
    \BlankLine
    
    \ForEach {horizontal split rule $s_1 \in S_1$}{
        Determine $\Delta I (s_1, X_1, Y)$\;
        \BlankLine
        \uIf{$\Delta I (s_1, X_1, Y) > \Delta I ^\star$}{
            $s^\star \gets s_1$ 
        }
    }
    \ForEach {vertical split rule $s_2 \in S_2$}{
        Determine $\Delta I (s_2, X_2, Y\T)$\;
        \BlankLine
        \uIf{$\Delta I (s_2, X_2, Y\T) > \Delta I ^\star$}{
            $s^\star \gets s_2$ 
        }
    }
    \BlankLine
    \Return $s^\star$
\end{function}


\begin{procedure}[tb]
    % Renamed from BatchAssignLeaves
    % AssignLeavesOxytrees -> AssignLeaves
    \caption{AssignLeaves($X_1$, $X_2$, Node): The proposed leaf assignment procedure for biclustering trees}
    \label{alg:batch assign leaves}
    
    \KwIn{A tree node and sets $X_1$ and $X_2$ of test instances in each dimension}
    \KwResult{A leaf is assigned to each interaction between $X_1$ and $X_2$}
    \SetKwFunction{self}{AssignLeavesOxytrees}
    \SetKwFunction{node}{Node}
    \SetKwFunction{leaf}{Leaf}
    \SetKwFunction{leftchild}{LeftChild}
    \SetKwFunction{rightchild}{RightChild}
    \BlankLine
    
    \uIf{\node is a \leaf}{
        Assign \leaf to all the interactions between $X_1$ and $X_2$\label{ln:final leaf assignment}\;
        \Return
    }
    \BlankLine
    
    Get the \leftchild and the \rightchild of \node\;
    Get the feature $f$ and threshold $t$ that define the split rule of \node\;
    \BlankLine
    
    \uIf{\node represents a horizontal split}{
        %Use the split rule of Node to partition $X_1$ into $X_\text{1, left}$ and $X_\text{1, right}$
        $X_\text{1, right}\gets \emptyset$\;
        $X_\text{1, left}\gets \emptyset$\;
        \ForEach(
            \tcp*[f]{Use the split rule to partition $X_1$}\label{ln:x segregation}
        ){
            $x_1 \in X_1$
        }{
            %\uIf{$\mathbf{s_1}(x_1)$}{
            \uIf{$x_1^f>t$}{
                $X_\text{1, right}\gets X_\text{1, right} \cup\{x_1\}$\;
            }
            \uElse{
                $X_\text{1, left}\gets X_\text{1, left} \cup\{x_1\}$\;
            }
        }
        \self($X_\text{1, left}$, $X_2$, \leftchild)\;
        \self($X_\text{1, right}$, $X_2$, \rightchild)\;
    }
    \uElse{
        %Use the split rule of Node to partition $X_1$ into $X_\text{1, left}$ and $X_\text{1, right}$
        Do the analogous procedure to partition $X_2$ into $X_\text{2, left}$ and $X_\text{2, right}$\;
        \self($X_1$, $X_\text{2, left}$, \leftchild)\;
        \self($X_1$, $X_\text{2, right}$, \rightchild)\;
    }
\end{procedure}



\FloatBarrier

\section{Proof of Oxytrees Compatibility Theorem}
\label{sec:proof}

% \begin{definition}[Dimension-non-increasing functions]
% \label{ass:well-behaved}
% For any considered function $f$ between two vector spaces, we require that the dimension of its image be not larger than the dimension of its domain:
% 
% \begin{equation}
%     \mathrm{dim}(\mathrm{Im}(f)) \le \mathrm{dim}(\mathrm{Dom}(f)).
% \end{equation}
% \end{definition}
% 
% Assumption \ref{ass:well-behaved} excludes, for example, space-filling curves~\cite{}, which are not relevant to our application setting.
% Smooth functions in general are contemplated, as a result of the Rank-Nullity Theorem~\cite{}.
% % [Given our applied setting]
% % [exclude space-filling curves]
% % [includes linear and smooth functions rank-nullity theorem]
% 
% \begin{definition}[Impurity function]
% Let $\Y \subseteq \R$ be the output space of the machine learning problem being modeled and let $\I \subseteq \R$ be called the impurity space. Let $I\colon \Y^{n\times m} \to \I$. $I$ is a well-defined impurity function iff
% [norm? triangle inequality?]
% 
% [define for bipartite learning, row and column-permutation invariant]
% 
% [is expressed in terms of dimension-non-incresing functions]
% 
% \end{definition}

\begin{definition}
Let $S_n$ be the set of all permutations on $n$ elements.
\end{definition}

\begin{definition}[Permutation invariance]
\label{def:permutation invariance}
A function $f:\R^n \to \R^m$ is permutation invariant iff, for all $P \in S_n$ and $\x \in \R^n$, $f(P(\x)) = f(\x)$.
% TODO: this is not how permutation invariance is defined in the zaheer paper though, why is it equivalent? this is important to be able to use theorem 10 in the current wording
\end{definition}

\begin{definition}
%Let $f:\R^n \to \R^m$.
Let $\sim_P$ be the equivalence relation
\begin{equation}
    \x \sim_P \x' \iff \exists P \in S_n : P(\x) = \x'
\end{equation}
\end{definition}

\begin{definition}
Let $[\x]_P$ be the equivalence class of $\x\in\R^n$ induced by $\sim_P$:
\begin{equation}
    \label{eq:equivalence class}
    [\x]_P=\{\x' \in \R^n: \x' \sim_P \x\}
\end{equation}
\end{definition}

\begin{definition}
Let $\R^n/\sim_P:=\{[\x]_P\}$ denote the set of all equivalence classes $[\x]_P$, i.e. the quotient space over $\R^n$ induced by $\sim_P$. 
%
\end{definition}

\begin{lemma}
\label{lemma:permutation invariance}
$f:\R^n \to \R^m$ is permutation invariant iff
%it can be factored.
there exists $b : \R^n/\sim_P \to \R^m$ such that $f = b \circ \pi$,
%where $\;\pi : \R^n \to \R^n/\sim_P$ is the canonical surjection $\x\mapsto [\x]_P$.
where $\;\pi$ is the canonical surjection $\x\mapsto [\x]_P$.
\end{lemma}

\begin{proof}
Let $P' \in S_n$ and $\x \in \R$. Now let $\x'=P'(\x)$. Since $P'$ exists, by definition $\x'\sim_P \x$ and $[\x]_P=[\x']_P$. Then 
\begin{equation}
    \label{eq:p1}
    \pi(P'(\x)) = [P'(\x)]_P
    = [\x']_P
    = [\x]_P
    = \pi(\x)
\end{equation}
To prove sufficiency, if $b$ exist and $f=b\circ \pi$, then 
\begin{align}
    f(\x) &= b(\pi(\x))\nonumber\\
    &= b([\x]_P)\nonumber\\
    &= b([P'(\x)]_P)\nonumber\\
    &= b(\pi(P'(\x)))\nonumber\\
    &= f(P'(\x)))
\end{align}
This is valid for all $\x \in \R^n$ and $P'\in S_n$, since both were taken at random. Therefore $f$ is permutation invariant.

To show necessity, let $f:\R^n \to \R^m$ be a permutation invariant function. Define $b([\x]_P)=f(\x)$. Let $\x,\x'\in \R^n$.
If $b([\x]_P) \neq b([\x']_P)$, $f(\x) \neq f(\x')$ by definition.
Since $f$ is permutation invariant, $\nexists P\in S_n : \x' = P(\x)$. Therefore, $\x'\nsim_P \x$, which implies $[\x']_P \neq [\x]_P$, showing that $b([\x]_P) \neq b([\x']_P) \implies [\x]_P \neq [\x']_P$. Since $\x$ and $\x'$ were taken at random, this shows that $b$ is a well-defined function.
\end{proof}


\begin{definition}[Row-proxy representation]
\label{def:proxy}


Let $I: \Y^{m\times n} \to \I$.
% TODO: what is mathcal Y, what is Y tilde. Why not R (given the previous definition)?
If it exists, a row-wise proxy representation of $I$ is defined as two real functions $\mu_r : \Y^{m} \to \M_r$ and $\rho_r \colon \M^n_r \to \I$, where $\M_r \subseteq \R^k$, such that:
%
\begin{enumerate}[label=\arabic*., ref=\thedefinition{} item \arabic*]
    %\item $I(Y) = \rho_r (\{\mu_r(Y^{i}) \;\forall \; i \in [1, 2, \dots, n]\})$, being $Y^{i}$ the $i$-th row of $Y$;
    % \forall Y \in \Y^{m \times n}: 
    \item \label{def:I from rho mu} $I(Y) = \rho_r \begin{pmatrix} \mu_r(Y^{1}), & \mu_r(Y^{2}), &\dots, & \mu_r(Y^{n})\end{pmatrix}$, with $Y^{i}$ the $i$-th row of $Y$;
    \item $\rho_r$ is permutation invariant;
    %\item $\Y \subseteq \R^{n \times m}$;
    %\item $\tilde {\Y} \subseteq \R$;
    \item\label{def:compression} The dimension of $\M_r$ is independent from the the number of columns of $Y \in \Y^{n \times m}$, i.e. %$\exists \;\mu_r \mid \Y^m \to \M_r\quad \textrm{and} \quad$
    $\exists k\in \N^\ast$, $\forall m \in \N^\ast: \M_r \subseteq \R^k$.
    % TODO: maybe better to work with dim here as well? shouldn't that be Y^m btw? so exists k such that for all m: k = dim(M) < dim(Y^m)
\end{enumerate}
\end{definition}
%
\Cref{def:compression} represents the key aspect that allows the reduction of computational complexity.
We define
% \begin{gather*}
%     \pmb{\mu}_r \colon \Y^{n \times m} \to \M^n_r:
%     \\
%     \pmb{\mu}_r(Y) = \begin{pmatrix} \mu_r(Y^{1}), & \mu_r(Y^{2}), &\dots, & \mu_r(Y^{n})\end{pmatrix}
% \end{gather*}
\begin{align*}
    \pmb{\mu}_r \colon \Y^{n \times m} & \to \M^n_r:
    \\
    Y & \mapsto \begin{pmatrix} \mu_r(Y^{1}), & \mu_r(Y^{2}), &\dots, & \mu_r(Y^{n})\end{pmatrix}
\end{align*}
%
in order to express \cref{def:I from rho mu} more concisely as $I = \rho_r \circ \pmb{\mu}_r$.

%Additionally, we require that $\mu_r$ respects $\mu_r\; \textrm{surjective} \implies \textrm{dim}(\textrm{Dom}(\mu_r)) \le $.

A column-wise proxy representation is defined by functions $\mu_c$ and $\rho_c$, respecting the analogous representations for $I(Y^T)$ and respecting the analogous constraints.
% TODO: elaborate

%From now on we will phrase the text in the rows direction, but all results are valid also in the columns direction.

%[do not enforce mu to be equal]

%[3 more general: exists m such that k < m]

%[mu(Yi) is the row proxy]

\begin{definition}[Surjectivity]
\[
    c\colon \mathcal{A} \to \mathcal{C} \text{ is surjective} \iff \forall C \in \mathcal{C},\exists A \in \mathcal{A} : c(A) = C
\]
\end{definition}

\begin{lemma}
\label{lemma:surj}

If the composition of two functions is surjective, the outer-most component function is also surjective:
\begin{equation}
  \begin{tikzcd}
    & \mathcal{B} \arrow{dr}{b} &
    \\
    \mathcal{A} \arrow{ur}{a} \arrow[two heads]{rr}{c} & & \mathcal{C}
  \end{tikzcd}
\end{equation}
% \[
%     \exists \text{surjective} c\colon \mathcal{A} \twoheadrightarrow \mathcal{C} \mid c = b\circ a \implies b \text{ is surjective}
% \]
\begin{multline}
    \forall a: \mathcal{A} \to \mathcal{B}, b: \mathcal{B} \to \mathcal{C}, c: \mathcal{A} \to \mathcal{C}:
    \\
    :(c = b \circ a) \land (c \text{ surjective}) \implies b \text{ surjective}
\end{multline}
\end{lemma}
% \begin{proof}
% If $c = b \circ a$ is surjective, then for all $A \in \mathcal{A}$ we can take $B = a(A)$ such that
% %Since $\forall \; A\in \mathcal{A}, \;\exists\; B \in \mathcal{B} \mid B = b(A)$,
% \begin{equation*}
% \forall\; C \in \mathcal{C},\; \exists \; B \in \mathcal{B} \mid b(B) = C
% \therefore b \; \text{is surjective}
% \end{equation*}
% \end{proof}

\begin{proof}
Let $C \in \mathcal{C}$. Since $c$ is surjective, there exists an $A \in \mathcal{A}$ such that $c(A) = C$. Now let $B:=a(A)$. Then 
\begin{equation*}
    C = c(A) = (b\circ a)(A) = b(a(A)) = b(B),
\end{equation*}
showing that $b$ is also surjective (since $C$ was taken at random).
\end{proof}


\begin{definition}
Let $Y^{i\ast}$ and $Y^{\ast j}$ denote, respectively, the $i$-th row and $j$-th column of $Y\in \Y^{n \times m}$.
\end{definition}

\begin{definition}%[Column-permutation invariance]
% TODO: Ahh I see from the commented part now. It's not clear from this definition that this is defining column permutation invariance due to the structure of "Let (...) and let (...)".
Let $P_c\in S_m$ be a permutation of columns of $Y \in \Y^{n \times m}$ and let $\PP_c:=\Y^{n \times m}/\sim_c$ be the quotient space of $\Y^{n \times m}$ induced by the equivalence relation $\sim_c$ satisfying
\begin{multline}
     Y \sim_c Y' \iff \exists \;P_c \in S_m :\\
     : Y' = \begin{bmatrix}
         P_c(Y^{1 \ast}), &
         P_c(Y^{2 \ast}), &
         \dots, &
         P_c(Y^{n \ast})
     \end{bmatrix}
\end{multline}
\end{definition}
% TODO: P_c has already been defined earlier in this definition as a single permutation. Did you mean to define the space of all permutations earlier? Or should this \exists P_c just be dropped?
% In the latter case, the relation would no longer be symmetrical though and thus no longer a equivalence relation
% Maybe you meant to write the quotient space as Y^{nxm} / \sim instead?


The existence of row-proxies implies existence of $I = \rho_r \circ \pmb{\mu}_c$.
The existence of column-proxies implies existence of $I = \pi \circ b$.
Without loss of generality, we assume $I$ to be surjective by setting its codomain to its image: $\I = \text{Im}(I)$.
%[The existence of column-proxies implies $I$ is column-permutation-invariant and thus $I = b \circ a$.]
%[Let $\PP_c = \Y^{n \times m}/P_c$.]
%
%[By definition of the canonical map we infer the existence of a surjective map $a\colon\Y^{n \times m} \to \PP_c$.]
\begin{equation}
  \begin{tikzcd}
    &
     \M_r^n
        %%\arrow[swap, dashed, shift right=1]{dd}{f}
        %\arrow[dashed, shift left=1, "g" near end]{dd}{}
        %\arrow[swap, dashleftarrow, shift right=1, "f" near end]{dd}{}
        \arrow[swap, dashleftarrow, "f" near start]{dd}{}
        \arrow{dr}{\rho_r}%
    &
    \\
    \Y^{n \times m}
        \arrow{ur}{\pmb{\mu}_r}
        %\arrow{ur}{M_r}
        \arrow[swap, two heads]{dr}{\pi}
        %\arrow[bend left=60]{rr}{I}
        \arrow[crossing over, "I" near end, two heads]{rr}{}
    &
    &
    \I
    \\
    &
    %\Y^{n \times m}/P_c
    \PP_c
        \arrow[swap]{ur}{b}
        %\arrow[swap, dashed, shift right=1]{uu}{g}
        %\arrow[swap, dashed, shift right=1, "g" near end]{uu}{}
    &
  \end{tikzcd}
\end{equation}


\begin{lemma}
\label{lemma:surj functions}
$b$ and $\rho_r$ are surjective.
%If they exist, $b$, $\rho_r$, and $g$ are surjective.
\end{lemma}
\begin{proof}
By applying Lemma \ref{lemma:surj} to the compositions
$I = b \circ \pi$,
$I =  \rho_r \circ \pmb{\mu}_r$,
% $a = g \circ \pmb{\mu}_r$,
%one concludes, respectively, that $b$ and $\rho_r$, and $g$ are surjective.
one concludes, respectively, that $b$ and $\rho_r$ are surjective.
\end{proof}


\begin{lemma}
\label{lemma:f exists}
If $\pmb{\mu}_r: \Y \to \M^n_r$ exists, then there also exists $f: \PP_c \to \M^n_r$.
% TODO: uniqueness?
\end{lemma}
\begin{proof}
% Choose an element $Y^\ast \in [Y]_c$ for each equivalence class $[Y]_c \in \PP_c$.
% Define $f([Y^\ast]_c):=h(Y^\ast)$. Let $Y, Y'\in \Y^{n \times m}$. Then $f([Y]_c)\neq f([Y']_c)\implies h(Y) \neq h(Y') \implies Y \neq Y'$. 
% Therefore $f$ is a well-defined function.

Choose a random element $Y$ in each equivalence class $[Y]_{c} \in \PP_c$.
% Call $Y$ the \emph{canonical element} of $[Y]_{c}$.
Let $g\colon \PP_c \to \Y$ be the function that maps each equivalence class to its chosen element: $g([Y]_c) \mapsto Y$.
Define $f([Y]_c):=\pmb{\mu}_r\circ g([Y]_c)$. Let $Y_a, Y_b\in \Y^{n \times m}$. Then
%
\begin{align*}
    &f([Y_a]_c)\neq f([Y_b]_c) \implies&&\\
    &\implies \pmb{\mu}_r\circ g([Y_a]_c) \neq \pmb{\mu}_r\circ g([Y_b]_c) &&\text{(definition of $f$)}\\
    &\implies \pmb{\mu}_r(g([Y_a]_c)) \neq \pmb{\mu}_r(g([Y_b]_c)) \\
    &\implies g([Y_a]_c) \neq g([Y_b]_c) && (\because \pmb{\mu}_r\text{ is well-defined})\\
    &\implies [Y_a]_c \neq [Y_b]_c. && (\because\text{ eq. classes are disjoint})
\end{align*}
%
Therefore $f$ is a well-defined function.
%To show uniqueness, take 
%% Let $[p] \subset \PP_c$ represent the equivalence class of a point $p\in\PP_c$, defined as $\forall p_A, p_B \in \PP_c$, $[p_A] = [p_B] \iff b(p_A) = b(b_B)$. The set of all equivalence classes is the quotient space $\PP_c/b$.
%% TODO: maybe it would be more simple to define this as [p] = {q in P_c | b(p) = b(q)}
%% 
%% [They are disjoint partitions]
%% 
%% Since $\rho_r$ is surjective, there exists a bijection $\beta : \I \to \PP_c/b$ (the inverse of the canonical map).
%% 
%% [Explicitly show?]
%% 
%% % \[
%% % \forall\; \tilde y \in \I \; \exists \; q_A \in \M_r^n \mid b(q_A) = \tilde y \implies 
%% % \]
%% 
%% We can choose an element $p_\ast \in [p_\ast]$ for each equivalence class $[p_\ast] \in \PP_c$. Let $h: \PP_c/b \to \PP_c$ be a map $\forall \;[p_\ast] \in \PP_c$, $h([p])\to p_\ast \;\forall\;p \in [p_\ast]$.
%% Because $h([p]) \in [p]$, by definition
%% \(
%%     h([p_A]) \ne h([p_B])
%%     \implies [p_A] \ne [p_B]
%% \).
%% % [Invoking the Axiom of Choice?]
%% % [We can pick an element $p^\ast$ in each equivalence class]
%% %$\forall [p]_b \exists p^\ast \in [p]_b \colon h(p) = p^\ast \forall p \in [p]_b$
%% 
%% The composition $g = h \circ \beta \circ \rho_r$ is a well-defined function $g\colon \M_r \to \PP_c$, i.e. it satisfies $\forall \; q_A, q_B \in \M_r,\; g(q_A) \neq g(q_B) \implies q_A \neq q_B$:
%% \begin{align*}
%%     g(q_A) \neq g(q_B)
%%         &\implies h([\beta(\rho_r(q_A))]) \ne h([\beta(\rho_r(q_B))])\\
%%         &\implies [\beta(\rho_r(q_A))] \ne [\beta(\rho_r(q_B))]\\
%%         &\implies \beta(\rho_r(q_A)) \ne \beta(\rho_r(q_B))\\
%%         &\implies \rho_r(q_A) \ne \rho_r(q_B)\\
%%         &\implies q_A \ne q_B
%% \end{align*}
%% 
%% [Since the diagram is symmetric, $f$ also exists.]
\end{proof}


% Then $\forall \tilde y \in \I$.
%m 
% Given $\rho_r$ and $b$ are surjective,
% there exists a bijection between the quotient spaces $\Y^{n \times m}/P/b$ and $\M_r/\rho_r$. Therefore, the 
%m 
% %$\forall \tilde{y} \in \I \exists [q]_{\rho_r} \subset \M_r, [p]_{\rho_r} \subset \Y^{n \times m}/P_c \colon $
%m 
% %$\forall \tilde{y_A}, \tilde{y_B} \in \I \exists [q]_{\rho_r} \subset \M_r, [p]_{\rho_r} \subset \Y^{n \times m}/P_c \colon $
% $\forall p_A, p_B \in \Y^{n \times m}/P \exists [q]_{\rho_r} \subset \M_r, [p]_{\rho_r} \subset \Y^{n \times m}/P_c \colon $




% \begin{lemma}
% $g$ is not dimension-non-increasing.
% \end{lemma}
% 
% \begin{proof}
% 
% Definition \ref{def:proxy}.\ref{def:compression} implies that, in general, $\mathrm{dim}(\M_r^n)<\mathrm{dim}(\Y^{n \times m})$.
% 
% [$\mathrm{dim}(\PP_c) = \mathrm{dim}(\Y^{n \times m})$]
% 
% Since $g$ is surjective (Lemma \ref{lemma:surj functions}), $\mathrm{Im}(g) = \PP_c$, and therefore $\mathrm{dim}(\mathrm{Im}(g))>\mathrm{dim}(\mathrm{Dom}(g))$. This implies that there does not exist a function $g\colon \M_r^n\to \PP_c$ that satisfies Assumption \ref{ass:well-behaved}.
% 
% \end{proof}
% 
% \begin{lemma}
% \label{lemma:f well-behaved}
% $f$ satisfies the dimension-non-increasing assumption (Assumption \ref{ass:well-behaved}).
% \end{lemma}
% 
% \begin{proof}
% 
% Definition \ref{def:proxy}.\ref{def:compression} implies that, in general, $\mathrm{dim}(\M_r^n)<\mathrm{dim}(\Y^{n \times m})$.
% 
% [$\mathrm{dim}(\PP_c) = \mathrm{dim}(\Y^{n \times m})$]
% 
% [$\mathrm{Im}(f)\subseteq \M_r^n \implies \mathrm{dim}(\mathrm{Im}(f)) \le \mathrm{dim}(\M_r^n)$]
% 
% Therefore $\mathrm{dim}(\mathrm{Im}(f)) < \mathrm{dim}(\mathrm{Dom}(f))$, satisfying assumption \ref{ass:well-behaved}.
% 
% \end{proof}

\begin{lemma}
\label{lemma:single row permutation}
$I$ is invariant under permutations of elements of a single row:
\begin{equation*}
    I \left(
    \begin{bmatrix}
        Y^{1\ast}\\
        Y^{2\ast}\\
        \vdots\\
        Y^{k\ast}\\
        \vdots\\
        Y^{n\ast}\\
    \end{bmatrix}   
    \right)
    =
    I \left(
    \begin{bmatrix}
        Y^{1\ast}\\
        Y^{2\ast}\\
        \vdots\\
        P_c(Y^{k\ast})\\
        \vdots\\
        Y^{n\ast}\\
    \end{bmatrix}   
    \right)
\end{equation*}

\end{lemma}
{\color{red}TODO: CORRECTION PENDING}

\begin{proof}
The composition
\begin{align*}
    f \circ \pi (Y') = f \circ \pi ( P_c (Y'))
    =
    f \circ \pi \left(
    \begin{bmatrix}
        P_c(Y^{1\ast})\\
        P_c(Y^{2\ast})\\
        \vdots\\
        P_c(Y^{k\ast})\\
        \vdots\\
        P_c(Y^{n\ast})\\
    \end{bmatrix}
    \right)
    \\
    \begin{bmatrix}
        P_c(Y^{1\ast})\\
        P_c(Y^{2\ast})\\
        \vdots\\
        P_c(Y^{k\ast})\\
        \vdots\\
        P_c(Y^{n\ast})\\
    \end{bmatrix}
\end{align*}
\begin{align*}
    f \circ \pi (Y') = \pmb{\mu}_r (Y_c)
    =
    \begin{bmatrix}
        \mu_r(Y_c^{1\ast})\\
        \mu_r(Y_c^{2\ast})\\
        \vdots\\
        \mu_r(Y_c^{k\ast})\\
        \vdots\\
        \mu_r(Y_c^{n\ast})\\
    \end{bmatrix}
    \\
    \begin{bmatrix}
        P_c(Y^{1\ast})\\
        P_c(Y^{2\ast})\\
        \vdots\\
        P_c(Y^{k\ast})\\
        \vdots\\
        P_c(Y^{n\ast})\\
    \end{bmatrix}
\end{align*}

\begin{align*}
    f_c \circ \pi_c (Y') = \pmb{\mu}_r (Y_c)\\
    f_r \circ \pi_r (Y') = \pmb{\mu}_c (Y_r)\\
    f_c \circ \pi_c (Y') = \pmb{\mu}_r (Y_c)\\
    f_r \circ \pi_r (Y') = \pmb{\mu}_c (Y_r)
\end{align*}
%Let $Y \in \Y^{n\times m}$ and $P' \in S_m$. 
%% Since $f$ exists (\cref{lemma:f exists}) and $\pmb{\mu}_r = f \circ \pi$, it follows from \cref{lemma:permutation invariance} that $\pmb{\mu}_r$ is invariant under column permutations. Therefore, $\pmb{\mu}_r(P'(Y)) = \pmb{\mu}_r(Y)$. 
%Let $k \in \{1, 2, \dots, n\}$.
%Since $\pmb{\mu}_r$ applies $\mu_r$ in isolation to each row,
%$\mu_r(P'(Y^{k\ast})=\mu_r(Y^{k\ast})$. Then $\mu_r$ is permutation invariant.
%
%Since $\mu_r$ is permutation invariant,
%\begin{equation}
%    \pmb{\mu}_r(Y)
%    = \begin{bmatrix}
%        \mu_r(Y^{1\ast})\\
%        \mu_r(Y^{2\ast})\\
%        \vdots\\
%        \mu_r(Y^{k\ast})\\
%        \vdots\\
%        \mu_r(Y^{n\ast})\\
%    \end{bmatrix}
%    = \begin{bmatrix}
%        \mu_r(Y^{1\ast})\\
%        \mu_r(Y^{2\ast})\\
%        \vdots\\
%        \mu_r(P'(Y^{k\ast}))\\
%        \vdots\\
%        \mu_r(Y^{n\ast})\\
%    \end{bmatrix}
%\end{equation}
%showing that $I$ does not change if we permute the elements of a single row of its input.
\end{proof}

\begin{lemma}
    \label{lemma:single col permutation}
    $I$ is invariant under permutations of elements of a single column.
\end{lemma}
\begin{proof}
Assuming \cref{def:proxy} also in the direction of columns, we can apply the same reasoning that leads to \cref{lemma:single row permutation} to conclude that $I$ is also unchanged if we permute the elements of a single column of its input.
\end{proof}

%\begin{lemma}
%    $I$ is invariant under vertical permutations
%\end{lemma}
%\begin{proof}
%[$\forall \;\mathbf{y} \in \Y^n,\;\mu_c(P(\mathbf{y}))=\mu_c(\mathbf{y})$]
%[$P= (Y^{kj}\; Y^{lj})$]
%\end{proof}

\begin{theorem}[\citet{zaheer_deep_2017}]
\label{the:deepsets}

% A map $\phi\colon\mathcal{A}\to \mathcal{B}$ between finite-dimension vector spaces $\mathcal{A}$ and $\mathcal{B}$ is permutation-invariant if and only if
% $\exists\; \rho,\,\mu\ \mid\forall\; A \in \mathcal{A},\;\phi(A) = \rho \left(\sum_{a \in A}\mu(a)\right)$.
Let $f\colon[0,\; 1]^M \to \mathbb{R}$ be a permutation-invariant continuous function iff it has the representation
\begin{equation}
    f\begin{pmatrix}
        x_1, & x_2, & \dots, & x_M 
    \end{pmatrix}
    = \rho\left(\sum_{m=1}^M\mu(x_m)\right)
\end{equation}
for some continuous outer and inner functions $\rho \colon \mathbb{R}^{M+1}\to \mathbb{R}$ and $\mu \colon \mathbb{R}\to \mathbb{R}^{M+1}$ respectively. The inner function $\phi$ is independent on the function $f$.

%$\exists\; \rho,\,\mu\ \mid\forall\; A \in \mathcal{A},\;\phi(A) = \rho \left(\sum_{a \in A}\mu(a)\right)$.
\end{theorem}

\begin{theorem}
\label{the:oxytrees}
There exists both a row-wise and a column-wise proxy representation for $I(Y)$ if and only if there exists $\rho$ and $\mu$ functions such that 
\begin{equation}
    I(Y)=\rho\left(\sum_{ij}\mu(Y^{ij})\right)
\end{equation}
\end{theorem}
\begin{proof}
To prove sufficiency, let 
\begin{align}
    \mu_r(z_1, \dots, z_m) &= \sum_j \mu(z_j)\\
    \rho_r(z_1, \dots, z_n) &= \rho\left(\sum_{i}z_i\right)\\
    \mu_c(z_1, \dots, z_n) &= \sum_i \mu(z_i)\\
    \rho_c(z_1, \dots, z_m) &= \rho\left(\sum_{j}z_j\right)
\end{align}
%
Given \cref{the:deepsets}, these functions follow the definition of proxies (\cref{def:proxy}).

To prove necessity, let $Y \in \Y^{n \times m}$. Let $(Y^{ij}\; Y^{kl})$ be the permutation of elements $Y^{ij}$ and $Y^{kl}$. $(Y^{ij}\; Y^{kl})$ can be written as a composition of permutations in a single column or row:
\begin{equation}
    (Y^{ij}\; Y^{kl}) = (Y^{ij}\; Y^{il}) \circ (Y^{il}\; Y^{jl}) \circ (Y^{ij}\; Y^{il})
\end{equation}
\Cref{lemma:single row permutation} implies that $I((Y^{ij}\; Y^{il})(Y)) = I(Y)$ and \cref{lemma:single col permutation} implies that $I((Y^{ij}\; Y^{kj})(Y)) = I(Y)$. Therefore
\begin{align}
I(&(Y^{ij}\; Y^{kl})(Y))=\nonumber\\
&=I((Y^{ij}\; Y^{il}) \circ (Y^{il}\; Y^{jl}) \circ (Y^{ij}\; Y^{il})(Y))\nonumber\\
&=I((Y^{ij}\; Y^{il}) \circ (Y^{il}\; Y^{jl})(Y))\nonumber\\
&=I((Y^{ij}\; Y^{il}) (Y))\nonumber\\
&=I(Y)
\end{align}
This is valid for any indices $i,j,k,l$, since they were taken at random.
Since any permutation of the elements of $Y$ can be written as a composition of permutations $(Y^{ij}\; Y^{kl})$, $I$ is invariant under any permutation of elements in $Y$. From \cref{the:deepsets}, it follows that $I$ can be expressed as
\begin{equation}
    I\begin{pmatrix}
        Y^{11}, & \dots, & Y^{nm}
    \end{pmatrix}
    = \rho\left(\sum_{j=1}^m\sum_{i=1}^n\mu(Y^{ij})\right)
\end{equation}
for some continuous $\rho \colon \mathbb{R}^{M+1}\to \mathbb{R}$ and $\mu \colon \mathbb{R}\to \mathbb{R}^{M+1}$.
% Let $\text{vec}(Y)$ be the vectorized $Y$, defined as $\text{vec}(Y)^{m \cdot i + j}:= Y^{ij}$.
% \begin{multline}
%     \begin{bmatrix}
%         Y^{ij} & Y^{il}\\
%         Y^{kj} & Y^{kl}
%     \end{bmatrix}
%     \to
%     \begin{bmatrix}
%         Y^{il} & Y^{ij}\\
%         Y^{kj} & Y^{kl}
%     \end{bmatrix}
%     \to
%     \\
%     \to
%     \begin{bmatrix}
%         Y^{il} & Y^{kl}\\
%         Y^{kj} & Y^{ij}
%     \end{bmatrix}
%     \to
%     \begin{bmatrix}
%         Y^{kl} & Y^{il}\\
%         Y^{kj} & Y^{ij}
%     \end{bmatrix}
% \end{multline}
\end{proof}

\begin{corollary}
The impurity followed by \citet{pliakos_global_2018} does not allow both row and proxy representations:
\begin{equation}
    \label{eq:pbct}
    I_\text{PBCT}(Y) = \frac{1}{m}\sum_j\left(\frac{1}{n}\sum_i (Y^{ij})^2 - \left(\frac{1}{n}\sum_i (Y^{ij})\right)^2\right)
\end{equation}
\end{corollary}
\begin{proof}
   From the representation in \cref{the:oxytrees}, we have that $I(Y) = I(\text{vec}(Y))$.
   Take a matrix $Y'$ for which columns are constant. From \cref{eq:pbct}, $I_\text{PBCT}(Y')=0 \neq I_\text{PBCT}(\text{vec}(Y'))$.
\end{proof}

%\section{Proof of Oxytrees Compatibility Theorem [Was not included]}
%\label{sec:proof}
%
%% \begin{equation}
%%     \exists \rho, \phi \; \colon \; I(Y) = \rho \left( \sum_i \phi(Y^i) \right) \;\forall \;Y \iff I = \tilde \rho \left( \sum_{ij} \tilde \phi(Y^{ij})\right)
%% \end{equation}
%
%\begin{definition}
%    Let $\M(\R)$ be the space of all real matrices of arbitrary dimensions.
%\end{definition}
%
%\begin{definition}
%    Let $P_\text{rows}(Y)$ and $P_\text{cols}(Y)$ be arbitrary permutations of rows and columns of
%    %a real matrix $Y$, respectively.
%    $Y\in \M(\R)$, respectively.
%\end{definition}
%%
%%Let $I(\cdot)$
%[TODO]
%
%\begin{definition}
%    Finally, let $I\colon \M(\R)\to \R$
%    be an impurity function satisfying $I(Y) = I(P_\text{rows}(\phi(P_\text{cols}(Y))))$.
%\end{definition}
%%
%
%An important result from \cite{zaheer_deep_2017}
%
%\begin{theorem}
%    \label{the:permutation invariance}
%    $f \colon \R^{m} \to \R^n$ is permutation invariant if and only if
%    $\exists\; g, $ such that $f(\y) = g\left(\sum_i h(\y^i)\right)$.
%\end{theorem}
%
%from \cref{the:permutation invariance} and \cref{def:unstructured}
%
%\begin{equation}
%    I(Y) = f_2 \left( \sum_j g_2\left( \sum_i h_2(Y^{ij}) \right) \right)
%\end{equation}
%
%from \cref{the:permutation invariance} and \cref{def:permutation invariance}
%
%\begin{equation}
%    I(Y) = f_1 \left( \sum_i \tilde g_1(Y^{i}) \right)
%\end{equation}
%
%% from \cref{the:permutation invariance} and 
%from \cref{def:unstructured} and \cref{def:permutation invariance}
%
%\begin{equation}
%    I(Y) = I(P(Y^T)^T)
%\end{equation}
%
%\begin{equation}
%    I(Y) = f_1 \left( \sum_i \tilde g_1(P(Y^{i})) \right)
%\end{equation}
%
%if $\tilde g_1$ commutes with $P$
%
%\begin{equation}
%    I(Y) = f_1 \left( \sum_i P(\tilde g_1(Y^{i})) \right)
%\end{equation}
%
%and $f_1$ must be permutation invariant, so that
%
%\begin{equation}
%    I(Y) = f_2 \left( \sum_j g_2\left( \sum_i \tilde g_1(Y^{i}) \right)\right)
%\end{equation}
%
%and no information is gained (already shown).
%
%if $\tilde g_1$ does not commute with $P$, $\tilde g_1$ must be permutation-invariant and
%
%\begin{equation}
%    I(Y) = f_1 \left( \sum_i g_1\left(\sum_j h_1(Y^{ij})\right) \right)
%\end{equation}
%
%for $\tilde g_1$ to commute with $P$, its domains must be compatible. \Cref{def:row-summarizability} defines that $k$ is independent from $m$ and therefore the commutation is not valid in general.
%
%\begin{equation}
%    \exists \rho \;\colon\; 
%    I(Y) = f_1 \left( \sum_j g_1\left( \sum_i h_1(Y^{ij}) \right) \right)
%\end{equation}
%
%\Cref{the:bgso text} is reformulated into
%
%\begin{multline}
%    \exists f_1, g_1, h_1, f_2, g_2, h_2
%    \\
%    \mid
%    f_2 \left( \sum_j g_2\left( \sum_i h_2(Y^{ij}) \right) \right)
%    = f_1 \left( \sum_i g_1\left( \sum_j h_1(Y^{ij}) \right) \right)
%    \;\forall \; Y \in \R^{*\times *}
%    \\
%    \iff g_1 \text{ or } g_2 \text{ is distributive.}
%\end{multline}
%
%\begin{multline}
%    \exists\\
%    f_1:\mathbb{G}_1^m \to \R \\
%    f_2:\mathbb{G}_2^n \to \R \\
%    g_1:\mathbb{H}_1^m \to \mathbb{G}_1 \\
%    g_2:\mathbb{H}_2^n \to \mathbb{G}_2 \\
%    h_1:\mathbb{Y} \to \mathbb{H}_1 \\
%    h_2:\mathbb{Y} \to \mathbb{H}_2
%    \\
%    \text{such that } g_2 \text{ is not injective}
%    % \mathbb{H}_2
%    \\
%    \text{and such that}
%    \\
%    f_2 \left( \sum_j g_2\left( \sum_i h_2(Y^{ij}) \right) \right)
%    = f_1 \left( \sum_i g_1\left( \sum_j h_1(Y^{ij}) \right) \right)
%    \;\forall \; Y \in \R^{*\times *}
%    \\
%    \iff g_1 \text{ or } g_2 \text{ is distributive.}
%\end{multline}
%
%% \begin{equation}
%%     f_2 \left( \sum_j g_2\left( \sum_i h_2(Y^{ij}) \right) \right)
%%     = f_1 \left( \sum_i g_1\left( \sum_j h_1(Y^{ij}) \right) \right)
%%     \iff g_2 \text{ is linear}
%% \end{equation}
%
%% \begin{equation}
%%     g_2 \text{ is linear}
%%     \implies
%%     \exists \rho \;\colon\; 
%%     % \rho \left( \sum_{ij} \phi(Y^{ij}) \right)
%%     \rho \left( \sum_{ij} \phi(Y^{ij})\right)
%%     = f_2 \left( \sum_j g_2\left( \sum_i h_2(Y^{ij}) \right) \right)
%% \end{equation}
%
%% set $\rho = g_2 \circ f_2$ and $h_2 = \phi$.
%
%% \begin{equation}
%%     g_2 \text{ is not linear}
%%     \implies
%%     %\nexists \rho \;\colon \;
%%     % \rho \left( \sum_{ij} \phi(Y^{ij})\right)
%%     \nexists \; f_1,g_1,h_1 \;\colon \;
%%     f_1 \left( \sum_i g_1\left( \sum_j h_1(Y^{ij}) \right) \right)
%%     = f_2 \left( \sum_j g_2\left( \sum_i h_2(Y^{ij}) \right) \right)
%%     \;\forall \;Y, f_2, g_2, h_2
%% \end{equation}
%% 
%% suppose exists
%
%% \begin{equation}
%%     \exists \; f_1,g_1,h_1 \;\colon \;
%%     f_1 \left( \sum_i g_1\left( \sum_j h_1(Y^{ij}) \right) \right)
%%     = f_2 \left( \sum_j g_2\left( \sum_i h_2(Y^{ij}) \right) \right)
%%     \;\forall \;Y, f_2, g_2, h_2
%% \end{equation}
%
%% \begin{equation}
%%     Y = \begin{bmatrix}
%%         a_1 & a_2 & \cdots & a_n \\
%%         \vdots & \vdots & \ddots & \vdots \\
%%         a_1 & a_2 & \cdots & a_n
%%     \end{bmatrix}
%% \begin{equation}
%
%
%
%%\begin{equation}
%%    Y = \begin{bmatrix}
%%        P_1(\y) & P_2(\y)& \dots& P_n(\y)
%%    \end{bmatrix}^T
%%\end{equation}
%\begin{equation}
%    \tilde Y^i = P_i(Y^i)
%\end{equation}
%
%\begin{equation}
%\end{equation}
%
%\begin{equation}
%    \sum_j g_2\left( \sum_i h_2(Y^{ij}) \right)
%    = f \left( \sum_i g_1\left( \sum_j h_1(Y^{ij}) \right) \right)
%\end{equation}
%
%\begin{equation}
%    \sum_i g_2\left( \sum_j h_2(Y^{ij}) \right)
%    = f \left( \sum_j g_1\left( \sum_i h_1(Y^{ij}) \right) \right)
%\end{equation}
%
%\begin{align}
%    \sum_j g_2\left( \sum_i h_2(Y^{ij}) \right)
%    = f \left( \sum_i g_1\left( \sum_j h_1(Y^{ij}) \right) \right)
%    \\
%    \sum e_m[g_2] \circ \sum_i e_{nm}[h_2]\circ P_j^i(Y)
%    = f \circ \sum e_m[g_1] \circ \sum_j e_{nm}[h_1]\circ P_j^i(Y)
%    \\
%    \sum e_m[g_2] \circ \sum_i P_j^i \circ e_{nm}[h_2](Y)
%    = f \circ \sum e_m[g_1] \circ \sum_j P_j^i\circ e_{nm}[h_1](Y)
%    \\
%    \sum e_m[g_2] \circ \sum_i P_j^i \circ e_{nm}[h_2](Y)
%    %= f \circ \sum e_m[g_1] \circ \sum_j e_{nm}[h_1](Y)
%    = \sum e_m[g_2] \circ \sum_i e_{nm}[h_2](Y)
%\end{align}
%
%\begin{align}
%    \sum e_m[g_2] \circ \sum_i P_i^j \circ e_{nm}[h_2](Y)
%    = f \circ \sum e_m[g_1] \circ \sum_j P_i^j\circ e_{nm}[h_1](Y)
%    \\
%    \sum e_m[g_2] \circ \sum_i e_{nm}[h_2](Y)
%    = f \circ \sum e_m[g_1] \circ \sum_j P_i^j\circ e_{nm}[h_1](Y)
%    \\
%    f \circ \sum e_m[g_1] \circ \sum_j e_{nm}[h_1](Y)
%    = f \circ \sum e_m[g_1] \circ \sum_j P_i^j\circ e_{nm}[h_1](Y)
%\end{align}
%\begin{align}
%    \sum e_m[g_2] \circ \sum_i P_i \circ e_{nm}[h_2](Y)
%    = f \circ \sum e_m[g_1] \circ \sum_j P_i\circ e_{nm}[h_1](Y)
%    \\
%    \sum e_m[g_2] \circ \sum_i e_{nm}[h_2](Y)
%    = f \circ \sum e_m[g_1] \circ \sum_j P_i\circ e_{nm}[h_1](Y)
%    \\
%    f \circ \sum e_m[g_1] \circ \sum_j e_{nm}[h_1](Y)
%    = f \circ \sum e_m[g_1] \circ P_i\circ \sum_j e_{nm}[h_1](Y)
%\end{align}
%\hrule
%
%\begin{equation}
%    Y = \begin{bmatrix}
%        a
%    \end{bmatrix}
%\end{equation}
%
%\begin{equation}
%    %\exists \; f_1,g_1,h_1 \;\colon \;
%    f_1 \left( g_1\left( h_1(a) \right) \right)
%    = f_2 \left( g_2\left( h_2(a) \right) \right)
%\end{equation}
%
%\begin{equation}
%    Y = \begin{bmatrix}
%        a_1 & a_2 & \cdots & a_n
%    \end{bmatrix}
%\end{equation}
%
%\begin{equation}
%    %\exists \; f_1,g_1,h_1 \;\colon \;
%    f_1 \left( g_1\left( \sum_j h_1(a^j) \right) \right)
%    = f_2 \left( \sum_j g_2\left( h_2(a^{j}) \right) \right)
%\end{equation}
%
%let $f = f_1^{-1} \circ f_2$
%
%\begin{equation}
%     g_1\left( \sum_j h_1(a^j) \right)
%    = f \left( \sum_j g_2\left( h_2(a^{j}) \right) \right)
%\end{equation}
%
%\begin{equation}
%    f \left( \sum_j g_2\left( \sum_i h_2(Y^{ij}) \right) \right)
%    = \sum_i g_1\left( \sum_j h_1(Y^{ij}) \right)
%\end{equation}
%
%\begin{equation}
%    f \left( \sum_j g_2\left( \sum_i h_2(Y^{ij}) \right) \right)
%    = \sum_i  f \left( \sum_j g_2\left( h_2(Y^{ij}) \right) \right)
%\end{equation}
%
%\begin{equation}
%    Y = \begin{bmatrix}
%        a_1 & a_2 & \cdots & a_n
%    \end{bmatrix}^T
%\end{equation}
%
%
%\begin{equation}
%    f \circ  g_2\left( \sum_i h_2(a^{i}) \right)
%    = \sum_i  f \circ  g_2\left( h_2(a^{i}) \right)
%\end{equation}
%
%
%\begin{equation}
%    f \circ  g_2 = l_2
%\end{equation}
%
%
%the same in the other direction
%
%\begin{equation}
%    Y = \begin{bmatrix}
%        a_1 & a_2 & \cdots & a_n
%    \end{bmatrix}^T
%\end{equation}
%
%\begin{equation}
%    %\exists \; f_1,g_1,h_1 \;\colon \;
%    f_1 \left( \sum_i g_1\left( h_1(a^i) \right) \right)
%    = f_2 \circ g_2\left( \sum_ih_2(a^{i}) \right)
%\end{equation}
%
%\begin{equation}
%    f^{-1} \left( \sum_i g_1\left( h_1(a^i) \right) \right)
%    = g_2\left( \sum_ih_2(a^{i}) \right)
%\end{equation}
%
%\begin{equation}
%    \sum_j g_2\left( \sum_i h_2(Y^{ij}) \right)
%    = f^{-1} \left( \sum_i g_1\left( \sum_j h_1(Y^{ij}) \right) \right)
%\end{equation}
%
%\begin{equation}
%    \sum_j f^{-1} \left( \sum_i g_1\left( h_1(Y^ij) \right) \right)
%    = f^{-1} \left( \sum_i g_1\left( \sum_j h_1(Y^{ij}) \right) \right)
%\end{equation}
%
%\begin{equation}
%    Y = \begin{bmatrix}
%        a_1 & a_2 & \cdots & a_n
%    \end{bmatrix}
%\end{equation}
%
%\begin{equation}
%    \sum_j f^{-1} \circ g_1\left( h_1(Y^ij) \right)
%    = f^{-1} \circ g_1\left( \sum_j h_1(Y^{ij}) \right)
%\end{equation}
%
%
%\begin{equation}
%    f^{-1} \circ  g_1 = l_1
%\end{equation}
%
%
%
%\begin{equation}
%    f \left( \sum_j g_2\left( \sum_i h_2(Y^{ij}) \right) \right)
%    = \sum_i g_1\left( \sum_j h_1(Y^{ij}) \right)
%\end{equation}
%
%\begin{equation}
%    Y = \begin{bmatrix}
%        a & a & \cdots & a \\
%        \vdots & \vdots & \ddots & \vdots \\
%        a & a & \cdots & a
%    \end{bmatrix}
%\end{equation}
%
%\begin{equation}
%    f \left( m \; g_2\left( n\; h_2(a) \right) \right)
%    = n\; g_1\left( m \; h_1(a) \right)
%\end{equation}
%
%\begin{equation}
%    Y = \begin{bmatrix}
%        a_1 & a_0 & \cdots & a_0 \\
%        a_0 & a_2 & \cdots & a_0 \\
%        \vdots & \vdots & \ddots & \vdots \\
%        a_0 & a_0 & \cdots & a_n
%    \end{bmatrix}
%\end{equation}
%
%$h_1 (a_0) = 0$
%
%\begin{equation}
%    f \circ g_2\circ h_2
%    = g_1 \circ h_1
%\end{equation}
%
%\begin{equation}
%    g_2 =  f^{-1} \circ l_2
%\end{equation}
%\begin{equation}
%    g_1 =  f \circ l_1
%\end{equation}
%
%\begin{equation}
%    l_2\circ h_2
%    = g_1 \circ h_1
%    \implies h_2
%    = l_2^{-1} \circ g_1 \circ h_1
%\end{equation}
%
%\begin{equation}
%    f \left( \sum_j f^{-1} \circ l_2 \left( \sum_i l_2^{-1} \circ g_1 \circ h_1(Y^{ij}) \right) \right)
%    = \sum_i g_1\left( \sum_j h_1(Y^{ij}) \right)
%\end{equation}
%
%\begin{equation}
%    f \left( \sum_j f^{-1} \circ l_2 \left( \sum_i l_2^{-1} \circ f \circ l_1 \circ h_1(Y^{ij}) \right) \right)
%    = \sum_i f \circ l_1\left( \sum_j h_1(Y^{ij}) \right)
%\end{equation}
%
%\begin{equation}
%    f \left( \sum_j f^{-1} \left( \sum_i f \circ l_1 \circ h_1(Y^{ij}) \right) \right)
%    = \sum_i f \left( \sum_j l_1 \circ h_1(Y^{ij}) \right)
%\end{equation}
%
%\begin{equation}
%    f \left( \sum_j f^{-1} \circ f \circ l_1 \circ h_1(a^{j}) \right)
%    = \sum_i f \circ l_1 \circ h_1(a^{i})
%\end{equation}
%
%\begin{equation}
%    f \left( \sum_i l_1 \circ h_1(a^{i}) \right)
%    = \sum_i f \circ l_1 \circ h_1(a^{i})
%\end{equation}
%
%
%$g_1$, $g_2$, and $f$ are all distributive.




\section{Asymptotic Computational Complexity of Biclustering Decision Trees}
\label{sec:complexity analysis}

In this appendix, we formally deduce the computational complexity expressions for the training and leaf assignment procedures of bipartite decision trees, comparing our method proposals to the previous algorithms introduced by \citet{pliakos_global_2018}.

\subsection{Training Time Complexity}
\label{sec:fit complexity analysis}

%TODO: fix n_1 or n notation
We assume that the numbers of instances in each dimension have the same order of complexity. Formally, let $n_1$ and $n_2$ be the number of rows in $X_1$ and $X_2$, respectively. We assume $n_1 \in \Theta(n_2)$.
%
From \ref{alg:find best split dt} alone, the complexity of the split search would be given by the number $|\s|$ of candidate splits times the time it takes to calculate $\Delta I$ for each of the splits: $\Theta(|\s| \;\Delta I(n_1))$. However, we must consider that it is often possible to obtain multiple values of $\Delta I$ for a single feature column $(X_a\T)^f$ much more efficiently using dynamic programming techniques. For instance, we can evaluate all split thresholds between two consecutive values of $(X_1\T)^f$ in $\Theta(\Delta I(n_1))$ time instead of $\Theta (n_1\;\Delta I (n_1))$. This is done by sorting $(X_1\T)^f$, applying the same permutation that sorted $(X_1\T)^f$ to the rows of $Y$, and iteratively obtaining $\tilde Y_1$ and $\tilde Y_2$ (\cref{eq:proxies general}) from which $\Delta I$ (\cref{eq:delta i}) can be calculated.
% (see Equations \ref{eq:i rows oxytrees} and \ref{eq:i rows pbct}).
 %
% 
Therefore, the final computational complexity depends on the number of features $m$
%of candidate features chosen to be evaluated (a hyperparameter),
and not on the total number of split candidates:
\begin{equation}
    \Theta(\text{\ref{alg:find best split dt}}) = \Theta(m \Delta I (n_1))
\end{equation}
%
In general, the complexity of $\Delta I$ is linear with respect to the number of elements of $Y$. Thus, for BICTR we have $\Delta I(n_1) \in \Theta (n_1^2)$, which results in
%
\begin{align}
    \label{eq:complexity find split gmo}
    \text{\ref{alg:find best split gmo}}
        &\in \Theta(m n_1^2)
    \text{.}
\end{align}
For Oxytrees, however, $\Delta I(n_1) \in \Theta (n_1)$, since we employ lower-dimensional proxies as described in \cref{sec:oxytrees training}. After introducing an $n_1^2$ factor to account for the time spent building the proxies, we have
\begin{align}
    \label{eq:complexity find split bgso}
    \text{\ref{alg:find best split bgso}}
        &\in \Theta(n_1^2 + m n_1)
    \text{.}
\end{align}

To estimate the complexity of the entire tree-building process, we consider the ideal scenario of a balanced decision tree.\footnote{Notice that the balanced scenario \emph{disfavors} our proposed optimization. Specifically, our method improves the complexity of the node-wise impurity calculation, which is dominated by the largest node in each level. A balanced tree results in the fastest decrease in the maximum node size per level.}
We also consider that the dimension of a split alternates at each level, so that a matrix $Y$ is separated into four equally-sized partitions after two levels. This results in the following recurrence relationship:
%
\begin{equation}
    T(n_1) = 4T\left(\frac{n_1}{2}\right) + \text{\ref{alg:find best split dt}}(n_1, m)
    \label{eq:recurrence}
\end{equation}
%
$T(n_1)$ is the time to build the tree from an $n_1$ by $n_1$ interaction matrix, and \ref{alg:find best split dt}$(n_1,\; m)$ in this case is the time taken to select a split.
The algorithm complexity of such recursive functions then follows the master theorem~\cite{cormen_introduction_2022}. The theorem gives the time complexity of a function $T$ obeying the recurrence relation $T(n) = aT(n/b) + F(n)$ and $c=\log_b a$:
%
\begin{equation}
    T(n) \in \begin{cases}
        \Theta(n^c) & \text{if } F(n) \in O(n^{c-\epsilon})\text{,} \\
        \Theta(n^c \log^{k+1} n) & \text{if } F(n) \in \Theta(n^c \log^k n)\text{,}\\
        \Theta(F(n)) & \text{if } F(n) \in \Omega(n^{c+\epsilon})
            \\&\quad\text{and }F(n)\text{ is regular.}
    \end{cases}
    \label{eq:master_theorem}
\end{equation}
%
In \cref{eq:master_theorem}, $\epsilon$ is a positive infinitesimal constant and $k$ is any non-negative integer. For the third case, we say a function $F(n)$ is regular if it satisfies the \emph{regularity condition}: $aF(n/b) \leq qF(n)$, for some constant $q < 1$ and all sufficiently large $n$~\cite{cormen_introduction_2022}.
% 
For our tree algorithms, \cref{eq:recurrence} shows that $c=2$ and $F(n)$ represents the split search function.
Applying \cref{eq:master_theorem} to \cref{eq:complexity find split bgso} and \cref{eq:complexity find split gmo} finally results in \cref{eq:complexity oxytree} and \cref{eq:complexity pbct}, respectively.


\subsection{Inference Time Complexity}
\label{sec:predict complexity analysis}
Let $|X|$ denote the number of rows in a matrix $X$.
Let $n_\text{train}$ represent the order of complexity of the number of training samples in both dimensions ($|X_\text{1, train}| \in \Theta(n_\text{train})$ and $|X_\text{2, train}| \in \Theta(n_\text{train})$). Similarly, $|X_\text{1, test}| \in \Theta(n_\text{test})$ and $|X_\text{2, test}| \in \Theta(n_\text{test})$.
For the traditional procedure of traversing the tree for each dyad individually, the complexity is given by
\begin{equation}
    \text{IndividualInference} \in \Theta(n_\text{test}^2\log\, n_\text{train})
\end{equation}
in which $n_\text{test}^2$ represents the total number of test dyads and $\log\, n_\text{train}$ represents the depth of the tree.
For the proposed batch inference procedure (\cref{sec:batch inference}), we first analyse the case in which $n_\text{test} > n_\text{train}$. The total time spent is given by
\begin{equation}
    \label{eq:time batch inference}
    \text{\ref{alg:batch assign leaves}}%(n_\text{train},\; n_\text{test})
    \in \Theta( T(n_\text{train}) + n_\text{test}^2)
    \text{,}
\end{equation}
%
in which the term $n_\text{test}^2$ represents the final leaf assignment for each test dyad (\cref{ln:final leaf assignment} of \cref{alg:batch assign leaves}) and $T$ represents the recursive procedure followed until the leaves. Similarly to the rationale behind \cref{eq:recurrence}, we consider a balanced tree in order to obtain
%
\begin{equation}
    T(n_\text{train}) = 4T\left(\frac{n_\text{train}}{2}\right) + n_\text{test}(n_\text{train})
    \text{.}
    \label{eq:recurrence predict}
\end{equation}
%
The last term corresponds to the linear process of segregating the instances in one of the dimensions (\cref{ln:x segregation} of \cref{alg:batch assign leaves}), considering that the expected number of test instances in a node will depend on the number of training instances in it.
Applying \cref{eq:master_theorem}, we see that the $n_\text{test}^2$ term of \cref{eq:time batch inference} dominates in all cases, no matter how $n_\text{test}$ asymptotically depends on $n_\text{train}$. That is,
\begin{equation}
    \label{eq:complexity batch inference n test greater}
    \text{\ref{alg:batch assign leaves}}
    %(n_\text{train},\; n_\text{test})
    \in \Omega(n_\text{test}^2)
\end{equation}
for $n_\text{test} > n_\text{train}$.
For the scenario in which $n_\text{test} < n_\text{train}$,
%
\begin{multline}
    \label{eq:time batch inference}
    \text{\ref{alg:batch assign leaves}}%(n_\text{train},\; n_\text{test}) 
    \in
        \\
        \in \Theta(
            T(n_\text{test}) 
            + 2 n_\text{test}^2 (\log\;n_\text{train} - \log\; n_\text{test}))
        )
    \text{,}
\end{multline}
%
with
%
\begin{equation}
    T(n_\text{test}) = 4T\left(\frac{n_\text{test}}{2}\right) + n_\text{test}
    \text{.}
    \label{eq:recurrence predict 2}
\end{equation}
%
$T$ corresponds again to the recursive procedure of \ref{alg:batch assign leaves}, this time considering that the number of test instances in each node approximately halves with each tree level, until there is a single test instance per node in a given level. The second term of \cref{eq:time batch inference} accounts for the remaining path that each of these instances must still traverse: $n_\text{test}^2$ is the total number of test instances and $(\log\;n^2_\text{train} - \log\; n^2_\text{test})$ is the expected length of the remaining path.
Since $n_\text{test} < n_\text{train}$, only two cases are investigated with \cref{eq:master_theorem}:
%
\begin{multline}
    \label{eq:complexity batch inforence n test lower}
    \text{\ref{alg:batch assign leaves}}%(n_\text{train},\; n_\text{test})
        \in\\
        \in \begin{cases}
            \Theta(n_\text{test}^2) & \text{if }n_\text{test} \in \Theta(n_\text{train})\text{,}\\
            \Theta(n_\text{test}^2 \log\;n_\text{train}) & \text{if }n_\text{test} \in o(n_\text{train})
            \text{.}
        \end{cases}
\end{multline}
%
Combining the conclusions of \cref{eq:complexity batch inference n test greater} and \cref{eq:complexity batch inforence n test lower}, we finally arrive at \cref{eq:complexity batch inference}.
\begin{align}
    \label{eq:complexity batch inference}
    %\text{\ref{alg:batch assign leaves}}(n_\text{train},\; n_\text{test})
    \text{\ref{alg:batch assign leaves}}
        \in \begin{cases}
            \Theta(n_\text{test}) \text{ if }n_\text{test} \in \Omega(n_\text{train})\text{,}\\
            \Theta(n_\text{test} \log(n_\text{train})) \text{ otherwise.}
        \end{cases}
\end{align}

\section{Details on Comparison Methods}
\label{sec:baselines details}

This section presents further details on the baseline methods used in this study, listed in \cref{sec:baselines}.

We included NRLMF \cite{liu_neighborhood_2016,liu_lpi-nrlmf_2017,liu_predicting_2020}, since it is, to the best of our knowledge, the most relevant matrix factorization method that can handle the fully inductive case.
It does so by regularizing the loss function to preserve neighborhoods between the original feature space and the latent space. Such neighborhoods can then be used to infer latent feature vectors (and thus predictions) for unseen instances.
NRLMF was successfully applied to several interaction prediction tasks, namely drug-target~\cite{liu_neighborhood_2016}, long non-coding RNA-protein~\cite{liu_lpi-nrlmf_2017}, and microRNA-lncRNA interactions~\cite{liu_predicting_2020}
%That is, it can predict new interactions without needing to reconstruct itself, making it applicable to all cases evaluated (Section X, EXPLANATION OF CROSS VALIDATION).

Oxytrees[Deep, YR] is included to validate our induction algorithm against BICTR~\cite{pliakos_drug-target_2020}. NRLMF is used by both methods to reconstruct the output space prior to training the forest.

Since logistic regression is a popular leaf model for model trees, we have included a baseline variant of our method using it with the GSO approach, Oxytrees[Logistic]. Furthermore, it also validates the use of RLS-Kron as the leaf model.   

Further, usual deep learning methods~\cite{huang_moltrans_2020,dehghan_tripletmultidti_2023} are not directly comparable since they are designed specifically for an application, e.g., drug-target \cite{liu2024ssldti}, RNA-disease \cite{tian2024mgcnss}. To account for neural networks, we included a multi-layer perceptron trained using the GSO approach as a baseline. 

As investigated in a previous work of ours~\cite{ilidio_fast_2024}, the semi-supervised variant of PBCT presents higher computational complexity, nonetheless it does not consistently lead to improvement in terms of performance \cite{ilidio_fast_2024}. Thus, we decide not to include it.       

The MLP architecture was selected between 10 layers of 100 neurons, 5 layers of 100 neurons, $\begin{bmatrix} 200 & 100 & 100 & 100 & 50 \end{bmatrix}$ and $\begin{bmatrix} 1024 & 512 & 256 & 128 & 64 & 32 \end{bmatrix}$, each number representing the size of each layer. The remaining hyperparameters were left as the default values of the Scikit-Learn package~\cite{pedregosa_scikit-learn_2011}, version 1.3.0.

Following the original WkNNIR study~\cite{liu_drug-target_2022}, the values evaluated for the number of neighbors used by WkNNIR were $\{1,2,3,5,7,9\}$, and for the learning rate were $\{0.1,0.2,\cdots,1\}$.

The Ridge regularization parameter of RLS-avg, RLS-Kron and BLMNII, was set to 1.
%
RLS-Kron, RLS-avg and BLMNII use a linear combination of a Gaussian network kernel (calculated from $Y$) with the original similarity kernel before training the models. We evaluated the values $\{0, 0.1, 0.25, 0.5, 0.75, 0.9, 1\}$ for the weight of this combination. The width of the Gaussian kernels was set to $1$, following the original authors.
%
The network kernel was not utilized for the RLS-Kron models in the leaves of Oxytrees.

% The Ridge regularization parameter of all linear models, including RLS-Kron, was set to 1, and no network-based kernel was used for .

NRLMF used a factor of $c=5$ to prioritize positive instances, as originally~\cite{liu_neighborhood_2016}. A maximum of 100 iterations with a minimum relative improvement of $10^{-5}$ for the loss function was performed. The number of neighbors was chosen between 3, 5 and 10, and the number of latent components was either 50 or 100. All remaining hyperparameters were randomly sampled from a log-uniform distribution between $2^{-4}$ and $2$. Due to the large hyperparameter space, 100 random combinations of hyperparameters were evaluated per outer fold for NRLMF, in contrast to the exhaustive search used with the other baselines. When used in combination with the forests, the hyperparameters were evaluated on the basis of the output of NRLMF, and not the final output of the forest. Furthermore, when combining NRLMF with forests, the original positive annotations are recovered, and only the NRLMF predictions for the negative (missing) annotations are kept~\cite{pliakos_drug-target_2020}.

All biclustering forests (BICTR and Oxytrees) utilized Extremely Randomized~\cite{geurts_extremely_2006} versions of the biclustering trees. That is, one random split candidate was selected for each candidate feature in each node. No bootstrapping of samples was performed, and all forests utilized 200 trees.
Following \cite{pliakos_network_2019,pliakos_drug-target_2020}, Oxytrees[Deep, YR] and BICTR had their trees grown until their maximum depth (reaching leaves with a single type of annotation) and evaluated $\lceil \sqrt{m}\rceil$ feature candidates in each node, from the $m$ total features. The remaining Oxytrees used minimum leaf dimensions of 5 by 5 as the only stopping criterion, and evaluated all $m$ features.

WkNNIR~\cite{liu_drug-target_2022} and BICTR~\cite{pliakos_global_2018} propose slight variations for the semi-inductive settings. We did not use them, since we primarily focus on inductive models.



\section{Additional Results}
\label{sec:additional results}
\FloatBarrier
%This section gathers additional figures and tables.
\begin{figure*}[tbh]
    \centering
    \begin{subfigure}{\textwidth}
        \includegraphics[width=.49\textwidth]{results/literature_methods/TT/all_datasets/critical_difference_diagrams/AUROC (Inductive).pdf}
        \hfill
        \includegraphics[width=.49\textwidth]{results/literature_methods/TT/all_datasets/critical_difference_diagrams/AUPRC (Inductive).pdf}
        \caption{Positives masking percentage = 0\% (Inductive setting)}
    \end{subfigure}
    
    \begin{subfigure}{\textwidth}
        \includegraphics[width=.49\textwidth]{results/literature_methods/TT_25/all_datasets/critical_difference_diagrams/AUROC (Inductive).pdf}
        \hfill
        \includegraphics[width=.49\textwidth]{results/literature_methods/TT_25/all_datasets/critical_difference_diagrams/AUPRC (Inductive).pdf}
        \caption{Positives masking percentage = 25\% (Inductive setting)}
    \end{subfigure}
    
    \begin{subfigure}{\textwidth}
        \includegraphics[width=.49\textwidth]{results/literature_methods/TT_50/all_datasets/critical_difference_diagrams/AUROC (Inductive).pdf}
        \hfill
        \includegraphics[width=.49\textwidth]{results/literature_methods/TT_50/all_datasets/critical_difference_diagrams/AUPRC (Inductive).pdf}
        \caption{Positives masking percentage = 50\% (Inductive setting)}
    \end{subfigure}
    
    \begin{subfigure}{\textwidth}
        \includegraphics[width=.49\textwidth]{results/literature_methods/TT_75/all_datasets/critical_difference_diagrams/AUROC (Inductive).pdf}
        \hfill
        \includegraphics[width=.49\textwidth]{results/literature_methods/TT_75/all_datasets/critical_difference_diagrams/AUPRC (Inductive).pdf}
        \caption{Positives masking percentage = 75\% (Inductive setting)}
    \end{subfigure}   

    \caption{Comparisons of AUROC (left) and AUPRC (right) of the proposed Oxytrees against previous methods for inductive bipartite learning. The analysis is conducted across four different positive interaction masking percentages.}
    \label{fig:literature inductive}
\end{figure*}


\begin{figure*}[tbh]
    \centering
    \begin{subfigure}{\textwidth}
        \includegraphics[width=.49\textwidth]{results/literature_methods/TT/all_datasets/critical_difference_diagrams/AUROC (Semi-inductive).pdf}
        \hfill
        \includegraphics[width=.49\textwidth]{results/literature_methods/TT/all_datasets/critical_difference_diagrams/AUPRC (Semi-inductive).pdf}
        \caption{Positives masking percentage = 0\% (Semi-inductive setting)}
    \end{subfigure}
    
    \begin{subfigure}{\textwidth}
        \includegraphics[width=.49\textwidth]{results/literature_methods/TT_25/all_datasets/critical_difference_diagrams/AUROC (Semi-inductive).pdf}
        \hfill
        \includegraphics[width=.49\textwidth]{results/literature_methods/TT_25/all_datasets/critical_difference_diagrams/AUPRC (Semi-inductive).pdf}
        \caption{Positives masking percentage = 25\% (Semi-inductive setting)}
    \end{subfigure}
    
    \begin{subfigure}{\textwidth}
        \includegraphics[width=.49\textwidth]{results/literature_methods/TT_50/all_datasets/critical_difference_diagrams/AUROC (Semi-inductive).pdf}
        \hfill
        \includegraphics[width=.49\textwidth]{results/literature_methods/TT_50/all_datasets/critical_difference_diagrams/AUPRC (Semi-inductive).pdf}
        \caption{Positives masking percentage = 50\% (Semi-inductive setting)}
    \end{subfigure}
    
    \begin{subfigure}{\textwidth}
        \includegraphics[width=.49\textwidth]{results/literature_methods/TT_75/all_datasets/critical_difference_diagrams/AUROC (Semi-inductive).pdf}
        \hfill
        \includegraphics[width=.49\textwidth]{results/literature_methods/TT_75/all_datasets/critical_difference_diagrams/AUPRC (Semi-inductive).pdf}
        \caption{Positives masking percentage = 75\% (Semi-inductive setting)}
    \end{subfigure}   

    \caption{Comparisons of AUROC (left) and AUPRC (right) of the proposed Oxytrees against previous methods for semi-inductive bipartite learning. The analysis is conducted across four different positive interaction masking percentages.}
    \label{fig:literature semi-inductive}
\end{figure*}

\begin{figure*}[tbh]
    \centering
    \begin{subfigure}{\textwidth}
        \includegraphics[width=.49\textwidth]{results/literature_methods/TT_25/all_datasets/critical_difference_diagrams/AUROC (Transductive).pdf}
        \hfill
        \includegraphics[width=.49\textwidth]{results/literature_methods/TT_25/all_datasets/critical_difference_diagrams/AUPRC (Transductive).pdf}
        \caption{Positives masking percentage = 25\% (Transductive setting)}
    \end{subfigure}
    
    \begin{subfigure}{\textwidth}
        \includegraphics[width=.49\textwidth]{results/literature_methods/TT_50/all_datasets/critical_difference_diagrams/AUROC (Transductive).pdf}
        \hfill
        \includegraphics[width=.49\textwidth]{results/literature_methods/TT_50/all_datasets/critical_difference_diagrams/AUPRC (Transductive).pdf}
        \caption{Positives masking percentage = 50\% (Transductive setting)}
    \end{subfigure}
    
    \begin{subfigure}{\textwidth}
        \includegraphics[width=.49\textwidth]{results/literature_methods/TT_75/all_datasets/critical_difference_diagrams/AUROC (Transductive).pdf}
        \hfill
        \includegraphics[width=.49\textwidth]{results/literature_methods/TT_75/all_datasets/critical_difference_diagrams/AUPRC (Transductive).pdf}
        \caption{Positives masking percentage = 75\% (Transductive setting)}
    \end{subfigure}   

    \caption{Comparisons of AUROC (left) and AUPRC (right) of the proposed Oxytrees against previous methods for transductive bipartite learning. The analysis is conducted across different positive interaction masking percentages.}
    \label{fig:literature transductive}
\end{figure*}



% \begin{figure*}
%     \begin{subfigure}{.49\textwidth}
%         \centering
%         \includegraphics[width=\textwidth]{results/y_reconstruction/TT/all_datasets/critical_difference_diagrams/AUROC (Inductive).pdf}
%         %\includegraphics[width=\textwidth]{results/y_reconstruction/TT/all_datasets/critical_difference_diagrams/AUROC (Inductive).pdf}
%         \caption{Positives masking percentage = 0\%}
%     \end{subfigure}
%     \begin{subfigure}{.49\textwidth}
%         \centering
%         \includegraphics[width=\textwidth]{results/y_reconstruction/TT_25/all_datasets/critical_difference_diagrams/AUROC (Inductive).pdf}
%         %\includegraphics[width=\textwidth]{results/y_reconstruction/TT_25/all_datasets/critical_difference_diagrams/AUROC (Inductive).pdf}
%         \caption{Positives masking percentage = 25\%}
%     \end{subfigure}
%     
%     \begin{subfigure}{.49\textwidth}
%         \centering
%         \includegraphics[width=\textwidth]{results/y_reconstruction/TT_50/all_datasets/critical_difference_diagrams/AUROC (Inductive).pdf}
%         %\includegraphics[width=\textwidth]{results/y_reconstruction/TT_50/all_datasets/critical_difference_diagrams/AUROC (Inductive).pdf}
%         \caption{Positives masking percentage = 50\%}
%     \end{subfigure}
%     \begin{subfigure}{.49\textwidth}
%         \centering
%         \includegraphics[width=\textwidth]{results/y_reconstruction/TT_75/all_datasets/critical_difference_diagrams/AUROC (Inductive).pdf}
%         %\includegraphics[width=\textwidth]{results/y_reconstruction/TT_75/all_datasets/critical_difference_diagrams/AUROC (Inductive).pdf}
%         \caption{Positives masking percentage = 75\%}
%     \end{subfigure}   
% 
%     \caption{AUROC comparison in the inductive case, of Oxytrees with and without a leaf model, as well as using and not using matrix completion with NRLMF. See \cref{sec:ablation}.}
%     \label{fig:ablation inductive auroc}
% \end{figure*}

\onecolumn

\begin{table*}[tb]
    \tiny
    %\label{tab:auprc}
    \centering
    %\onecolumn
    \input{figures/results/latex_tables/0/AUPRC (Inductive)}
    \label{tab:auprc}
\end{table*}

\begin{table*}[tb]
    \tiny
    \centering
    \input{figures/results/latex_tables/0/AUROC (Inductive)}
    \label{tab:auroc}
\end{table*}


\begin{table*}[tb]
    \tiny
    \centering
    \input{figures/results/latex_tables/0/AUPRC (Transductive)}
    \label{tab:auprc transductive}
\end{table*}
\begin{table*}[tb]
    \tiny
    \centering
    \input{figures/results/latex_tables/0/AUROC (Transductive)}
    \label{tab:auroc transductive}
\end{table*}


\begin{table*}[tb]
    \tiny
    \centering
    \input{figures/results/latex_tables/0/AUPRC (LT)}
    \label{tab:auprc lt}
\end{table*}
\begin{table*}[tb]
    \tiny
    \centering
    \input{figures/results/latex_tables/0/AUROC (LT)}
    \label{tab:auroc lt}
\end{table*}


\begin{table*}[tb]
    \tiny
    \centering
    \input{figures/results/latex_tables/0/AUPRC (TL)}
    \label{tab:auprc tl}
\end{table*}
\begin{table*}[tb]
    \tiny
    \centering
    \input{figures/results/latex_tables/0/AUROC (TL)}
    \label{tab:auroc tl}
\end{table*}

\begin{figure*}
    \centering
    \begin{subfigure}{\textwidth}
        \includegraphics[width=.47\linewidth]{figures/results/latex_tables/0/AUROC (Inductive).pdf}
        %\caption{Training time as a function of the number of instances.}
        %\label{fig:pmp auroc inductive}
        \hfill
        \includegraphics[width=.47\linewidth]{figures/results/latex_tables/0/AUPRC (Inductive).pdf}
        %\caption{Inference time as a function of the number of instances.}
        %\label{fig:pmp auprc inductive}
        \caption{Performance drop in the \emph{inductive} case, as a function of PMP.}
    \end{subfigure}
    
    \begin{subfigure}{\textwidth}
        \includegraphics[width=.47\linewidth]{figures/results/latex_tables/0/AUROC (Semi-inductive).pdf}
        %\caption{Training time as a function of the number of instances.}
        %\label{fig:pmp auroc inductive}
        \hfill
        \includegraphics[width=.47\linewidth]{figures/results/latex_tables/0/AUPRC (Semi-inductive).pdf}
        %\caption{Inference time as a function of the number of instances.}
        %\label{fig:pmp auprc inductive}
        \caption{Performance drop in the \emph{semi-inductive} case, as a function of PMP.}
    \end{subfigure}
    \caption{AUROC (left) and AUPRC (right) for different positive masking percentages (PMP) divided by the corresponding value in the PMP = 0\% setting. The relative values are obtained for each fold. They are averaged across folds of each dataset and then across datasets, for all datasets in \cref{tab:datasets} and all the compared models of \cref{sec:baselines}. For the semi-inductive case, we considered scores on both the TL and LT sets (\cref{fig:train test split}) when averaging within each dataset.}
    \label{fig:pmp drop}
\end{figure*}

\begin{figure*}[tbp]
    \includegraphics[width=.31\textwidth]{figures/results/empirical_complexity_plots/empirical_complexity_predict.pdf}
    \hfill
    \includegraphics[width=.31\textwidth]{figures/results/empirical_complexity_plots/fit_time_Oxytree_BICTR.pdf}
    \hfill
    \includegraphics[width=.31\textwidth]{figures/results/empirical_complexity_plots/predict_time_Oxytree_BICTR.pdf}
    \caption{Extra results for the empirical complexity analysis.
    Single trees were applied to artificial datasets of different dimensions (\cref{sec:empirical complexity}).
    On the left, prediction times for the same trees used in \cref{fig:empirical complexity} (left).
    We used the last 10\% of the points of each curve to approximate the asymptotic complexity as $\Theta (n^\alpha)$. $\alpha$ ($\pm$ standard dev.) is estimated as
    %the slope in the log-log space.
    the slope of the linear regression in the log-log space.
    In the middle, training times of Oxytrees against BICTR and corresponding least squares linear regression (using the last 10\% of the points in \cref{fig:empirical complexity}).
    On the right, prediction times of Oxytrees against BICTR and corresponding least squares linear regression (using the last 10\% of the points of the figure in the left).
    The last points are used to approximate the asymptotic regime as well as possible.
    %BICTR leaf assignment terminated earlier due to memory limitations.
    }
 
    \label{fig:complexity extra}
\end{figure*}

\begin{figure*}
    \includegraphics[width=0.245\linewidth]{figures/results/ablation_relative/Inductive_AUPRC_00.pdf}
    \hfill
    \includegraphics[width=0.245\linewidth]{figures/results/ablation_relative/Inductive_AUPRC_25.pdf}
    \hfill
    \includegraphics[width=0.245\linewidth]{figures/results/ablation_relative/Inductive_AUPRC_50.pdf}
    \hfill
    \includegraphics[width=0.245\linewidth]{figures/results/ablation_relative/Inductive_AUPRC_75.pdf}
    
    \includegraphics[width=0.245\linewidth]{figures/results/ablation_relative/Semi-inductive_AUPRC_00.pdf}
    \hfill
    \includegraphics[width=0.245\linewidth]{figures/results/ablation_relative/Semi-inductive_AUPRC_25.pdf}
    \hfill
    \includegraphics[width=0.245\linewidth]{figures/results/ablation_relative/Semi-inductive_AUPRC_50.pdf}
    \hfill
    \includegraphics[width=0.245\linewidth]{figures/results/ablation_relative/Semi-inductive_AUPRC_75.pdf}
    
    \hfill
    \includegraphics[width=0.245\linewidth]{figures/results/ablation_relative/Transductive_AUPRC_25.pdf}
    \hspace{0.001\linewidth}
    \includegraphics[width=0.245\linewidth]{figures/results/ablation_relative/Transductive_AUPRC_50.pdf}
    \hspace{0.001\linewidth}
    \includegraphics[width=0.245\linewidth]{figures/results/ablation_relative/Transductive_AUPRC_75.pdf}
    
    \caption{
    AUPRC comparison of ensembles of Oxytrees using different components (\cref{tab:ablation}). To compare across datasets, the scores are divided by the score of the main model (Oxytrees). Each point represents a cross-validation fold, and the mean value is presented in the white boxes.
    Asterisks indicate significance in comparison to Oxytrees, averaging folds for each dataset ($\ast:p<0.05$, $\ast\ast:p<0.01$, $\ast\ast\ast:p<0.001$, Wilcoxon signed-rank).
    AUROC results are presented in \cref{fig:ablation extra auroc}.
    }
    \label{fig:ablation extra}
\end{figure*}

\begin{figure*}
    \includegraphics[width=0.245\linewidth]{figures/results/ablation_relative/Inductive_AUROC_00.pdf}
    \hfill
    \includegraphics[width=0.245\linewidth]{figures/results/ablation_relative/Inductive_AUROC_25.pdf}
    \hfill
    \includegraphics[width=0.245\linewidth]{figures/results/ablation_relative/Inductive_AUROC_50.pdf}
    \hfill
    \includegraphics[width=0.245\linewidth]{figures/results/ablation_relative/Inductive_AUROC_75.pdf}
    
    \includegraphics[width=0.245\linewidth]{figures/results/ablation_relative/Semi-inductive_AUROC_00.pdf}
    \hfill
    \includegraphics[width=0.245\linewidth]{figures/results/ablation_relative/Semi-inductive_AUROC_25.pdf}
    \hfill
    \includegraphics[width=0.245\linewidth]{figures/results/ablation_relative/Semi-inductive_AUROC_50.pdf}
    \hfill
    \includegraphics[width=0.245\linewidth]{figures/results/ablation_relative/Semi-inductive_AUROC_75.pdf}
    
    \hfill
    \includegraphics[width=0.245\linewidth]{figures/results/ablation_relative/Transductive_AUROC_25.pdf}
    \hspace{0.001\linewidth}
    \includegraphics[width=0.245\linewidth]{figures/results/ablation_relative/Transductive_AUROC_50.pdf}
    \hspace{0.001\linewidth}
    \includegraphics[width=0.245\linewidth]{figures/results/ablation_relative/Transductive_AUROC_75.pdf}
    
    \caption{
    AUROC comparison of ensembles of Oxytrees using different components (\cref{tab:ablation}). To compare across datasets, the scores are divided by the score of the main model (Oxytrees). Each point represents a cross-validation fold, and the mean value is presented in the white boxes.
    Asterisks indicate significance in comparison to Oxytrees, averaging folds for each dataset ($\ast:p<0.05$, $\ast\ast:p<0.01$, $\ast\ast\ast:p<0.001$, Wilcoxon signed-rank).
    }
    \label{fig:ablation extra auroc}
\end{figure*}

\begin{figure*}
    %\includegraphics[width=.48\linewidth]{figures/results/number_of_trees/Inductive_AUROC_0_BICTR_Oxytrees.pdf}
    %\hfill
    %\includegraphics[width=.495\linewidth]{figures/results/number_of_trees/Inductive_AUPRC_0_BICTR_Oxytrees.pdf}
    \includegraphics[width=0.245\linewidth]{figures/results/number_of_trees/Inductive_AUPRC_0_BICTR_Oxytrees.pdf}
    \hfill
    \includegraphics[width=0.245\linewidth]{figures/results/number_of_trees/Inductive_AUPRC_25_BICTR_Oxytrees.pdf}
    \hfill
    \includegraphics[width=0.245\linewidth]{figures/results/number_of_trees/Inductive_AUPRC_50_BICTR_Oxytrees.pdf}
    \hfill
    \includegraphics[width=0.245\linewidth]{figures/results/number_of_trees/Inductive_AUPRC_75_BICTR_Oxytrees.pdf}

    \includegraphics[width=0.245\linewidth]{figures/results/number_of_trees/Semi-inductive_AUPRC_0_BICTR_Oxytrees.pdf}
    \hfill
    \includegraphics[width=0.245\linewidth]{figures/results/number_of_trees/Semi-inductive_AUPRC_25_BICTR_Oxytrees.pdf}
    \hfill
    \includegraphics[width=0.245\linewidth]{figures/results/number_of_trees/Semi-inductive_AUPRC_50_BICTR_Oxytrees.pdf}
    \hfill
    \includegraphics[width=0.245\linewidth]{figures/results/number_of_trees/Semi-inductive_AUPRC_75_BICTR_Oxytrees.pdf}

    \includegraphics[width=0.245\linewidth]{figures/results/number_of_trees/Inductive_AUROC_0_BICTR_Oxytrees.pdf}
    \hfill
    \includegraphics[width=0.245\linewidth]{figures/results/number_of_trees/Inductive_AUROC_25_BICTR_Oxytrees.pdf}
    \hfill
    \includegraphics[width=0.245\linewidth]{figures/results/number_of_trees/Inductive_AUROC_50_BICTR_Oxytrees.pdf}
    \hfill
    \includegraphics[width=0.245\linewidth]{figures/results/number_of_trees/Inductive_AUROC_75_BICTR_Oxytrees.pdf}
   
    \includegraphics[width=0.245\linewidth]{figures/results/number_of_trees/Semi-inductive_AUROC_0_BICTR_Oxytrees.pdf}
    \hfill
    \includegraphics[width=0.245\linewidth]{figures/results/number_of_trees/Semi-inductive_AUROC_25_BICTR_Oxytrees.pdf}
    \hfill
    \includegraphics[width=0.245\linewidth]{figures/results/number_of_trees/Semi-inductive_AUROC_50_BICTR_Oxytrees.pdf}
    \hfill
    \includegraphics[width=0.245\linewidth]{figures/results/number_of_trees/Semi-inductive_AUROC_75_BICTR_Oxytrees.pdf}

%    \hfill
%    \includegraphics[width=0.245\linewidth]{figures/results/number_of_trees/Transductive_AUPRC_0_BICTR_Oxytrees.pdf}
%    \hspace{0.001\linewidth}
%    \includegraphics[width=0.245\linewidth]{figures/results/number_of_trees/Transductive_AUPRC_25_BICTR_Oxytrees.pdf}
%    \hspace{0.001\linewidth}
%    \includegraphics[width=0.245\linewidth]{figures/results/number_of_trees/Transductive_AUPRC_50_BICTR_Oxytrees.pdf}
%    \hspace{0.001\linewidth}
%    \includegraphics[width=0.245\linewidth]{figures/results/number_of_trees/Transductive_AUPRC_75_BICTR_Oxytrees.pdf}
 
    \caption{Comparison between Oxytrees and BICTR in terms of the expected number of trees to achieve 98\% of the score of 200 trees. Each fold is plotted, with different styles for each dataset.
    This quantity is estimated from the piece-wise linear interpolation in a plot of the score metric versus number of trees.
    LT and TL points are plotted separately in the semi-inductive plots (hence double the number or points).
    Comparisons of Oxytrees against Oxytrees without the leaf model (Oxytrees[Mean] in \cref{tab:ablation}) are shown in \cref{fig:number of trees no leaf model}
    }
%   The p-value for the comparison
%   (Wilcoxon signed-rank)
%   is displayed in the legend, based on the averaged scores per dataset.
    \label{fig:number of trees extra}
\end{figure*}

\begin{figure*}
    %\includegraphics[width=.48\linewidth]{figures/results/number_of_trees/Inductive_AUROC_0_BICTR_Oxytrees.pdf}
    %\hfill
    %\includegraphics[width=.495\linewidth]{figures/results/number_of_trees/Inductive_AUPRC_0_BICTR_Oxytrees.pdf}
    \includegraphics[width=0.245\linewidth]{figures/results/number_of_trees/Inductive_AUPRC_0_Oxytrees_Oxytrees[no leaf model].pdf}
    \hfill
    \includegraphics[width=0.245\linewidth]{figures/results/number_of_trees/Inductive_AUPRC_25_Oxytrees_Oxytrees[no leaf model].pdf}
    \hfill
    \includegraphics[width=0.245\linewidth]{figures/results/number_of_trees/Inductive_AUPRC_50_Oxytrees_Oxytrees[no leaf model].pdf}
    \hfill
    \includegraphics[width=0.245\linewidth]{figures/results/number_of_trees/Inductive_AUPRC_75_Oxytrees_Oxytrees[no leaf model].pdf}

    \includegraphics[width=0.245\linewidth]{figures/results/number_of_trees/Semi-inductive_AUPRC_0_Oxytrees_Oxytrees[no leaf model].pdf}
    \hfill
    \includegraphics[width=0.245\linewidth]{figures/results/number_of_trees/Semi-inductive_AUPRC_25_Oxytrees_Oxytrees[no leaf model].pdf}
    \hfill
    \includegraphics[width=0.245\linewidth]{figures/results/number_of_trees/Semi-inductive_AUPRC_50_Oxytrees_Oxytrees[no leaf model].pdf}
    \hfill
    \includegraphics[width=0.245\linewidth]{figures/results/number_of_trees/Semi-inductive_AUPRC_75_Oxytrees_Oxytrees[no leaf model].pdf}

    \includegraphics[width=0.245\linewidth]{figures/results/number_of_trees/Inductive_AUROC_0_Oxytrees_Oxytrees[no leaf model].pdf}
    \hfill
    \includegraphics[width=0.245\linewidth]{figures/results/number_of_trees/Inductive_AUROC_25_Oxytrees_Oxytrees[no leaf model].pdf}
    \hfill
    \includegraphics[width=0.245\linewidth]{figures/results/number_of_trees/Inductive_AUROC_50_Oxytrees_Oxytrees[no leaf model].pdf}
    \hfill
    \includegraphics[width=0.245\linewidth]{figures/results/number_of_trees/Inductive_AUROC_75_Oxytrees_Oxytrees[no leaf model].pdf}
   
    \includegraphics[width=0.245\linewidth]{figures/results/number_of_trees/Semi-inductive_AUROC_0_Oxytrees_Oxytrees[no leaf model].pdf}
    \hfill
    \includegraphics[width=0.245\linewidth]{figures/results/number_of_trees/Semi-inductive_AUROC_25_Oxytrees_Oxytrees[no leaf model].pdf}
    \hfill
    \includegraphics[width=0.245\linewidth]{figures/results/number_of_trees/Semi-inductive_AUROC_50_Oxytrees_Oxytrees[no leaf model].pdf}
    \hfill
    \includegraphics[width=0.245\linewidth]{figures/results/number_of_trees/Semi-inductive_AUROC_75_Oxytrees_Oxytrees[no leaf model].pdf}

%    \hfill
%    \includegraphics[width=0.245\linewidth]{figures/results/number_of_trees/Transductive_AUPRC_0_BICTR_Oxytrees.pdf}
%    \hspace{0.001\linewidth}
%    \includegraphics[width=0.245\linewidth]{figures/results/number_of_trees/Transductive_AUPRC_25_BICTR_Oxytrees.pdf}
%    \hspace{0.001\linewidth}
%    \includegraphics[width=0.245\linewidth]{figures/results/number_of_trees/Transductive_AUPRC_50_BICTR_Oxytrees.pdf}
%    \hspace{0.001\linewidth}
%    \includegraphics[width=0.245\linewidth]{figures/results/number_of_trees/Transductive_AUPRC_75_BICTR_Oxytrees.pdf}
 
    \caption{
    Comparison between Oxytrees and Oxytrees without the leaf model (Oxytrees[no leaf model] is Oxytrees[Mean] in \cref{tab:ablation})
    in terms of the expected number of trees to achieve 98\% of the score of 200 trees. This quantity is estimated from the piece-wise linear interpolation in a plot of the score metric versus number of trees.
    Each fold is plotted, with different styles for each dataset. LT and TL points are plotted separately in the semi-inductive plots (hence double the number or points).}
    \label{fig:number of trees no leaf model}
\end{figure*}

\FloatBarrier



\end{appendices}

\fi

\end{document}